% Version 1.1.1, 2026-10-15 (for spherical-ramsey v1.1.1; changes listed in the sections
% "Changes in version 1.1.1" and "Changes in version 1.1" and in CHANGELOG.md of the release).
% Every claim carries exactly one verification tag (see "Status conventions"); open problems are
% collected in the public conjecture ledger that accompanies this note (ledger/ in the release).
% Vocabulary: "Comparator-checked" only for statements with a fresh-runner CI receipt;
% "Lean-checked" only where the full release gate (CI, non-builder fidelity review, fresh-agent
% reproduction) has passed for that rope. Theorem C is published (doi:10.5281/zenodo.23339523): it
% is stated and cited here, not re-proved. The quotation of OpenAI's Theorem 1.1 is copied from the
% published note of [MR]. Pure ASCII source. No \TODO or \PH markers remain.
\documentclass[11pt]{amsart}
\usepackage[margin=1in]{geometry}
\usepackage{amsmath,amssymb,amsthm,mathrsfs}
\usepackage{mathtools}% psmallmatrix
\usepackage{booktabs}
\usepackage{xcolor}
\usepackage[hidelinks]{hyperref}
\setlength{\emergencystretch}{3em}% fewer overfull lines; no content change

% Robust, so that amsart's upper-cased title and headings leave the colour names alone.
\DeclareRobustCommand{\TODO}[1]{\textcolor{red}{[TODO: #1]}}
\DeclareRobustCommand{\PH}[1]{\textcolor{blue}{[#1]}}
\newcommand{\st}[1]{\textsf{[#1]}}
\newcommand{\Kriz}{K\v{r}\'{\i}\v{z}}
\newcommand{\Q}{\mathbb{Q}}
\newcommand{\R}{\mathbb{R}}
\newcommand{\C}{\mathbb{C}}
\newcommand{\HH}{\mathbb{H}}
\newcommand{\Qbar}{\overline{\mathbb{Q}}}
\newcommand{\Cl}{\mathrm{Cl}}
\newcommand{\Rea}{\operatorname{Re}}
\newcommand{\circmap}{\mathrm{circ}}
\newcommand{\BB}{\mathscr{B}}
\DeclareMathOperator{\PO}{PO}
\DeclareMathOperator{\Ort}{O}
\DeclareMathOperator{\CO}{CO}
\DeclareMathOperator{\PGL}{PGL}
\DeclareMathOperator{\PSL}{PSL}
\DeclareMathOperator{\SO}{SO}
\DeclareMathOperator{\Mat}{Mat}
\DeclareMathOperator{\Mob}{M\ddot{o}b}
\DeclareMathOperator{\adj}{adj}

\theoremstyle{plain}
\newtheorem{theorem}{Theorem}[section]
\newtheorem{proposition}[theorem]{Proposition}
\newtheorem{corollary}[theorem]{Corollary}
\newtheorem{lemma}[theorem]{Lemma}
\newtheorem{conjecture}{Conjecture}
\newtheorem*{oaithm}{Theorem 1.1}
\theoremstyle{definition}
\newtheorem{definition}[theorem]{Definition}
\newtheorem{remark}[theorem]{Remark}

\title[Two M\"obius groups suffice]{Two M\"obius groups suffice: Euclidean Ramsey sets on spheres of
every dimension, after the tensor classification}
\author{Will Strinz}
\date{Version 1.1.1, 15 October 2026}

\begin{document}

\begin{abstract}
OpenAI's classification of finite Euclidean Ramsey configurations \cite{OAI172} decides whether a
finite spherical set is Ramsey by a tensor condition over its coordinate field. We record
consequences of it; all of them rest on \cite{OAI172}, and the deep work is OpenAI's.
The main one is a \emph{two-group theorem}: for every $d\ge2$, a union of two M\"obius images of
finite subsets of $S^{d-1}$ with real algebraic coordinates is Ramsey, with no bound on the number
of points. The certificate pairs two Clifford-algebra graph forms, built from Vahlen matrices in a
normal form whose entries lie in the Lipschitz group. For $d\le6$ ($S^1$ to $S^5$) every identity
it needs is verified as an exact polynomial identity and the passage from identities to the
theorem is written out; for general $d$ the argument is a written proof, with its identity families
cross-checked on exact instances through $d=8$. Both levels have been reviewed only by AI agents.
Two groups are not the only mechanism. Off the circle this already follows from OpenAI's
quadric-independence criterion \cite[Prop.~7.3]{OAI172} (for example, seven points of $S^2$ with
algebraically independent coordinates are Ramsey and are not two-group); on the circle that criterion
reaches at most five points. On the circle we give exact certificates for $7$- and
$8$-point Ramsey sets that are not unions of two M\"obius images of algebraic sets, among them the
seven points $0,\theta,\dots,6\theta$ for transcendental $e^{i\theta}$; a rank criterion for
one-parameter cyclic sets of at most eight points, from which the eight points
$0,\theta,\dots,7\theta$ are Ramsey for every $\theta$; and a rigidity theorem: if the coordinate
field is regular over $\Q$ and five of the points have rational cross-ratios, Ramsey-ness is
equivalent to the two-group property. These are written proofs reviewed only by AI agents, except
the certificates, which are exact identities.
We also state, with their verification status: M\"obius invariance of Ramsey-ness, published
separately with a Lean proof \cite{MR}; the nine points $0,\theta,\dots,8\theta$ on a circle are
Ramsey if and only if $e^{i\theta}$ is algebraic (at $\theta=e$ this is equivalent to a
transcendence question that is, to our knowledge, open); algebraic circle points plus at most
three arbitrary points are Ramsey; and, on the non-Ramsey side, a phase diagram for rational plus
independent points on a circle and exact certificates for the smallest such configurations on
$S^2$ and $S^3$. The theorems on six points, on algebraic points plus three, and on the radian $9$-gon are proved
in Lean~4 on top of OpenAI's formalization and checked by Comparator on a fresh CI runner. We describe a second-order
necessary test, with computational evidence for its sufficiency, and state conjectures, including an
expected-dimension law for cyclic sets and measurable forms of Graham's conjecture. We observe that
Graham's conjecture holds trivially for colourings definable in an o-minimal structure (folklore
tameness), and, by a theorem of Dol'nikov and Karasev, for colourings that factor through boundedly
many polynomials. Several of the easier statements are
prior work by others (\S\ref{sec:credit}); version~1.1 corrects omissions in that credit
(see ``Changes in version~1.1'' below). A public ledger of the open problems, with replayable
starting states, accompanies the note.
\end{abstract}

\maketitle

\section*{Changes in version 1.1.1}

\begin{itemize}
\item \textbf{Written proofs of first-order rigidity and of the finite-type reformulation.} Version~1.1.0
labelled ledger entries SR-C34 and SR-C33(a)--(c) as proved \st{HAND} although only a summary of the
argument was published. The complete proofs are now in \S\ref{sec:rigidity}: Theorem~\ref{thm:rigidity}
(SR-C34), Proposition~\ref{prop:finitetype} (SR-C33(a)) and Proposition~\ref{prop:corankone}
(SR-C33(b),(c)), with the lemmas they use (the local ring at the diagonal, the obstruction, a Smith form,
and the duality between the 2-jet system and its dual space). Status: \st{HAND}: written proof in
\S\ref{sec:rigidity} (added in v1.1.1), independently reviewed twice by non-author AI agents (the first
review reconstructed the argument; a second, independent review found every step sound;
\S\ref{sec:disclosure}). By-products, proved there: the ``symmetric part'' conditions of
the 2-jet dual are automatic (Lemma~\ref{lem:Pi}), the dual form of the 2-jet test holds at every quadric
rank (Corollary~\ref{cor:jetdual}), and the rank criterion holds without its regularity hypothesis
(Corollary~\ref{cor:rigidcases}(c)).
\item \textbf{Scope of Theorem~C's Lean status.} The summary table now says what \S\ref{sec:mobius} says:
\st{LEAN} for the unit sphere and the maps $\varphi_L$; other spheres by a similarity, by hand.
\end{itemize}

\section*{Changes in version 1.1}

Version~1.0.0 overlooked a result of OpenAI and three appendix theorems of P\'alv\"olgyi that bear on
it; version~1.1 credits them and corrects what they change. The full list is in \texttt{CHANGELOG.md}
of the release.
\begin{itemize}
\item \textbf{OpenAI's quadric-independence criterion} \cite[Prop.~7.3]{OAI172} (Lean:
\texttt{spherical\_\allowbreak{}quadraticIndependent\_\allowbreak{}ramsey\_\allowbreak{}all}) is now credited
(\S\ref{sec:credit}, Remark~\ref{rem:quadric}). It implies that up to $(d+1)(d+2)/2-1$ points of
$S^{d-1}$ with algebraically independent coordinates are Ramsey, so non-two-group Ramsey sets off the
circle are prior work; Theorem~\ref{thm:NTG} is a statement about the circle, where that criterion
reaches at most five points. The ledger's remark that a certificate for seven generic points on $S^2$
would be ``the first certified Ramsey set beyond the two-group mechanism'' was wrong and is withdrawn.
\item \textbf{P\'alv\"olgyi's Theorems~A.11--A.13} \cite{Palvolgyi} are now credited: A.12 is the
non-Ramsey half of Conjecture~\ref{conj:C3} for $d\ge3$, and the ten-point cell of
Theorem~\ref{thm:NR} is its case $d=3$.
\item The scalar-certificate evidence of \S\ref{sec:jet} is re-scoped: on the tested cells it is
automatic. The searches of \S\ref{sec:credit} are stated as done (the version~1.0.0 text still carried
a marker for them), Theorem~C is registered on the Palomar registry, and two references on Vahlen
matrices are added.
\item \textbf{Addition: first-order rigidity} (ledger SR-C34). \st{HAND}: written proof in \S\ref{sec:rigidity} (added in
version~1.1.1), independently reviewed twice by non-author AI agents (the first review reconstructed the
argument; a second, independent review found every step sound); scope
\texttt{ANY}. Let $A$ lie on the unit sphere, let $D_1,\dots,D_k$ be a basis of the
derivations of its coordinate field, and let $\Delta(A)$ be the dual space of the 2-jet system of
\S\ref{sec:jet} without the trace condition. If $\Delta(A)=0$, then $A$ is Ramsey. The proof is a Smith
form over the local ring of $F\otimes_\Q F$ at the diagonal, the duality behind
Proposition~\ref{prop:J2}, and Nakayama's lemma. It contains OpenAI's quadric-independence criterion
(no quadric relation) and Theorem~\ref{thm:rank} (circle, one parameter), there without the regularity
hypothesis. A literature sweep (2026-10-15) found no prior statement of it. Its order-zero case is
OpenAI's Proposition~7.3, and its cases with algebraic coordinates, or with at most $m+4$ points in
affine dimension $m$, are prior work \cite{JDJones}; the new content is positive transcendence degree
together with quadric relations.
\item \textbf{Addition: failures without a scalar certificate} (ledger SR-C32, proved for four seeded
families of twelve and thirteen points on $S^2$; \st{CERT} plus glue, reviewed by one non-author AI
agent). They fail the 2-jet test and have no scalar (single-relation) derivation certificate, even over
$\C$. A matrix-weighted form of Lemma~\ref{lem:A1} still colours them (ledger SR-C35, \st{SKETCH},
not reviewed); it is an explicit derivation-energy instance of the colouring that OpenAI use for the
necessity half of their criterion \cite[Prop.~2.4]{OAI172}, and P\'alv\"olgyi's Theorem~A.4 is its
scalar-weight analogue.
\item Further new ledger entries (SR-C29 to SR-C35) and a frontier document accompany the release; each
entry states its own review status.
\end{itemize}

\noindent\textbf{Status conventions.} Every result carries one tag.
\st{LEAN}: proved in Lean~4 on top of OpenAI's unmodified Lean proof of \cite{OAI172}, with
\texttt{\#print axioms} reporting only \texttt{propext}, \texttt{Classical.choice},
\texttt{Quot.sound} (Lean v4.34.1 and Mathlib \texttt{d13f23b7} for the ropes of
Appendix~\ref{app:lean}; Lean v4.35.0-rc2 and Mathlib \texttt{065356127b} for the published package
\cite{MR}).
\st{CERT}: an exact symbolic identity (sympy, over $\Q$) plus written glue that is not
machine-checked.
\st{HAND}: a complete written argument, not machine-checked.
\st{SKETCH}: an outline.
\st{PROBE}: computational evidence, often modulo a prime; never a theorem.
For \st{LEAN} results we say \emph{Comparator-checked} when a stock Comparator \cite{Comparator}
run on a fresh CI runner accepted the statement against OpenAI's challenge definitions and refused
negative controls (the run number is given), and \emph{Lean-checked} when, in addition, a
non-builder fidelity review and a fresh-agent reproduction have passed; today that is
Theorems~\ref{thm:A}, \ref{thm:Bb} and \ref{thm:D} and, for the published package \cite{MR}, Theorem~\ref{thm:C}. All reviews cited below were performed by AI agents, unless stated
otherwise (\S\ref{sec:disclosure}). Each result also names the type of its coordinates over $\Q$
(\texttt{ALGEBRAIC}, \texttt{GENERIC\_OF}, \texttt{TRANSCENDENTAL\_PARAM},
\texttt{MOBIUS\_ORBIT\_OF}, \texttt{ANY}). ``Real'' alone is never a scope, because the verdict
can change between algebraic and transcendental members of one real family
(Theorem~\ref{thm:D}). CI run numbers below for the ropes are fresh-runner runs of this release's
own workflows (which re-run every rope and every fast exact script) on a private staging copy of the
release repository, on a tree identical to the released one (the release repository runs the same
workflows on every push); the runs for \cite{MR} are that package's. Script paths
(\texttt{certificates/\dots}, \texttt{probes/\dots}, \texttt{ropes/\dots}) are relative to the release
repository \url{https://github.com/wstrinz/spherical-ramsey}.

\section{Introduction}

A finite set $X$ in a Euclidean space is \emph{Ramsey} if for every $r\ge2$ there is $D$ such that
every $r$-colouring of $\R^D$ is constant on some congruent copy of $X$. Erd\H{o}s, Graham,
Montgomery, Rothschild, Spencer and Straus \cite{EGMRSS} showed that Ramsey sets are spherical;
Graham conjectured the converse \cite{Graham94}. Known Ramsey families include triangles
\cite{FR86}, simplices \cite{FR90}, trapezoids \cite{Kriz92}, sets with suitable transitive
symmetry \cite{Kriz91}, and regular polytopes \cite{Cantwell}. Leader, Russell and Walters
\cite{LRW12,LRW11} proposed a description in terms of transitive sets; their Conjecture~3 is
refuted in \cite[Cor.~7.5]{OAI172}. In September 2026 P\'alv\"olgyi \cite{Palvolgyi} refuted
Graham's conjecture with a seven-point non-Ramsey set on a circle, and shortly afterwards OpenAI
\cite{OAI172} classified the finite Ramsey configurations exactly, with a complete Lean proof.

The classification turns ``is $A$ Ramsey?'' into ``does a linear system over $F\otimes_\Q F$ have a
solution?'', where $F$ is the coordinate field of $A$. For algebraic $F$ the answer is always yes
(Theorem~\ref{thm:A}, prior work). For transcendental $F$ it depends on the configuration, and
nine generic points on a circle already fail \cite[Prop.~7.6]{OAI172}. The question this note
pursues is which transcendental configurations still pass. The M\"obius group gives one answer:
Ramsey-ness is M\"obius invariant (Theorem~\ref{thm:C}, published \cite{MR}), so every M\"obius
image of an algebraic set is Ramsey, whatever its transcendence degree. Our main theorem says that
\emph{two} such groups can be combined (Theorem~\ref{thm:C2}): the union of two M\"obius images of
algebraic sets, under unrelated maps, is Ramsey on spheres of every dimension. The certificate is
heuristically the polar form of a quadratic ``spinor to Lorentz matrix'' map, so it has exactly two
slots (Remark~\ref{rem:whytwo}), and three groups fail a second-order necessary test in every case we tried (\st{PROBE},
\S\ref{sec:jet}). Theorem~\ref{thm:D} shows that the boundary between the two regimes can sit on
a transcendence problem that is, to our knowledge, open.

Two groups are not the whole story (\S\ref{sec:beyond}). Off the circle this already follows from
OpenAI's quadric-independence criterion \cite[Prop.~7.3]{OAI172}: seven to nine points of $S^2$ with
algebraically independent coordinates are Ramsey and not two-group (Remark~\ref{rem:quadric}). On the
circle that criterion reaches at most five points, but there are Ramsey sets of
seven and eight points, with a single transcendental parameter, that are not two-group
(Theorem~\ref{thm:NTG}). For one-parameter cyclic sets of at most eight points a rank condition on
an explicit $n\times8$ matrix implies Ramsey (Theorem~\ref{thm:rank}); it gives the radian
$8$-gon for every $\theta$ (Corollary~\ref{cor:radian8}), one point short of Theorem~\ref{thm:D}.
In the other direction, five points with rational cross-ratios make the two-group property
necessary, for regular coordinate fields (Theorem~\ref{thm:rigid}). These results fit an
expected-dimension count for cyclic sets (Conjecture~\ref{conj:count}).

\medskip
{\small
\noindent\begin{tabular}{@{}p{0.13\textwidth}p{0.40\textwidth}p{0.27\textwidth}p{0.14\textwidth}@{}}
\toprule
Result & Statement & Status & Scope \\
\midrule
Thm~\ref{thm:C2} & two M\"obius images of algebraic sets on $S^{d-1}$ are Ramsey, every $d\ge2$
 & \st{CERT}+glue for $d\le6$; \st{HAND} for all $d$; AI-reviewed & M\"obius orbits \\
Thm~\ref{thm:C} & M\"obius invariance; published \cite{MR} & \st{LEAN} for the unit sphere (maps $\varphi_L$), Lean-checked \cite{MR}; other spheres by a similarity, \st{HAND} & any \\
Thm~\ref{thm:A} & algebraic coordinates $\Rightarrow$ Ramsey (prior work) & \st{LEAN}, Lean-checked & algebraic \\
Thm~\ref{thm:Ba} & $\le6$ points on a circle (prior work) & \st{LEAN}, Comparator-checked & any \\
Thm~\ref{thm:Bb} & algebraic circle points $+\le3$ arbitrary & \st{LEAN}, Lean-checked & alg.\ $\cup$ 3 any \\
Thm~\ref{thm:D} & radian 9-gon Ramsey iff $e^{i\theta}$ algebraic & \st{LEAN}, Lean-checked & transc.\ / alg. \\
Rem.~\ref{rem:quadric} & up to $(d+1)(d+2)/2-1$ generic points on $S^{d-1}$ (OpenAI's criterion; prior) & \st{LEAN} (OpenAI) + \st{HAND} glue, AI-reviewed & generic \\
Thm~\ref{thm:NTG} & non-two-group Ramsey sets of 7 and 8 circle points & \st{CERT}+glue & transc.\ param. \\
Thm~\ref{thm:rank} & rank criterion, one parameter, $n\le8$ & \st{HAND}, AI-reviewed & trdeg 1 \\
Cor~\ref{cor:radian8} & radian 8-gon Ramsey for every $\theta$ & \st{HAND}, AI-reviewed, exact rank & any \\
Thm~\ref{thm:rigid} & rigid-group absorption on the circle & \st{HAND}, AI-reviewed & regular field \\
Prop~\ref{prop:PD} & circle phase diagram & \st{HAND}, AI-reviewed & generic \\
Thm~\ref{thm:NR} & minimal non-Ramsey cells on $S^2$, $S^3$ & \st{CERT}+glue, AI-reviewed & generic \\
Prop~\ref{prop:J2} & the 2-jet test is necessary & \st{HAND}, AI-reviewed & any \\
Thm~\ref{thm:rigidity} & vanishing 2-jet dual space $\Rightarrow$ Ramsey (first-order rigidity) & \st{HAND}, reviewed twice (AI) & any \\
Conj.~\ref{conj:C1} & the 2-jet test is sufficient & evidence \st{PROBE} & any \\
Rem.~\ref{rem:definable} & Graham for o-minimal and polynomial-factor colourings & \st{HAND}, AI-reviewed & any \\
\bottomrule
\end{tabular}}

\subsection{Prior work and credit}\label{sec:credit}

\begin{itemize}
\item \textbf{OpenAI \cite{OAI172}} is the foundation: Theorem~1.1 and Proposition~2.1
(\S\ref{sec:prelim}), the examples ($\le5$ circle points Ramsey, Cor.~7.4; Leader--Russell--Walters
kites with transcendental parameter Ramsey, Cor.~7.5; nine generic circle points and a 12-point
three-squares set non-Ramsey, Props.~7.6 and 7.7), and the Lean proof on which all our Lean results
build. Version~1.0.0 of this note did not cite OpenAI's Proposition~7.3, the
\emph{quadric-independence criterion}: a spherical set is Ramsey if the vectors $p_ip_i^T$ are
linearly independent over $F$ (Lean: \texttt{spherical\_\allowbreak{}quadraticIndependent\_\allowbreak{}ramsey\_\allowbreak{}all}, Comparator
challenge \texttt{EuclideanRamseyQuadratic}; Cor.~7.4 is its circle case). It is vendored, unmodified,
in every rope of this release. Its consequences for generic points on spheres
(Remark~\ref{rem:quadric}) are prior work, not ours.
\item \textbf{P\'alv\"olgyi \cite{Palvolgyi}} gave the seven-point non-Ramsey heptagon. By the arXiv
comments and the disclosure in the paper itself, the whole manuscript was written by ChatGPT; its
appendix, which he labels as unchecked, consists of further results by ChatGPT that he has not
verified. Several appendix items bear on this note, all at the level of his derivation
certificates: Lemma~A.1, a general non-Ramsey criterion from a derivation and real weights, which we
re-derive and use in \S\ref{sec:nonramsey}; Theorem~A.3, its circle form, of which our exact
instance J2-F7 is a special case; A.2, that no certificate exists for at most six circle
points; Lemma~A.5, that the existence of a certificate is invariant under real M\"obius maps, with
the explicit statement that it ``does not assert that the Ramsey property itself is M\"obius
invariant''; Theorem~A.7, which has a M\"obius-invariant hypothesis without stating invariance of
Ramsey-ness, and which with A.3 gives the non-Ramsey half of Proposition~\ref{prop:PD}; and Theorem~A.10, the certificate-level analogue of Theorem~\ref{thm:Bb} and of
\cite[Cor.~1(c)]{MR} on the circle (a set that a M\"obius map sends to algebraic points, apart from at
most three, admits no certificate), with the note that this does not decide Ramsey-ness.
Three further appendix theorems were not cited in version~1.0.0: Theorem~A.11, that a fixed weighted
derivation-energy identity with all weights nonzero needs more than $2d+2$ points of an affinely
$d$-dimensional spherical set (so none exists on at most six points); Theorem~A.12, that for $d\ge3$ a set
of $(d+1)(d+2)/2$ points on $S^{d-1}$ whose stereographic parameters are algebraically independent is
not Ramsey, which is the non-Ramsey half of Conjecture~\ref{conj:C3} for $d\ge3$ and contains the
ten-point cell of Theorem~\ref{thm:NR}; and Theorem~A.13, explicit non-Ramsey spherical sets of $2d+3$
points in every affine dimension $d\ge2$.
\item \textbf{Theorem~\ref{thm:A} is prior work.} JD-Jones-ASES/six-concyclic-lean \cite{JDJones}
(root commit \texttt{0703d7d}, 2026-10-07) proved in Lean, at the level of OpenAI's field
criterion, that algebraic coordinates satisfy the criterion, that any six concyclic points do (by the
same two-triple certificate we use for Theorem~\ref{thm:Ba}), and that at most $m+4$ points with
$m$-dimensional affine span do; its Ramsey conclusions are conditional on a hypothesis
\texttt{hsuff} standing for OpenAI's sufficiency direction, and its paper text argues the sharpness
($d+5$) that makes seven the minimum on a circle. Rogerhu12/openai-math-extensions \cite{Rogerhu}
(commit \texttt{f8a2bb2}, 2026-10-07) proves Theorem~\ref{thm:A} in \TeX{} with the same
separability-idempotent argument. eliottcassidy2000/math \cite{Eliott} (2026-10-07) restates it in
an audit note, conditionally on \cite{OAI172}. atimics/math2 \cite{Atimics} (commit \texttt{8ab51fa},
2026-10-08) proved it end to end in Lean on OpenAI's proof, for algebraic configurations with full
affine span (its Lean statement carries a spanning hypothesis); its written derivation extends this
to sets similar to algebraic ones, a reduction not formalized at that commit. Our Lean proofs of
Theorems~\ref{thm:A} and \ref{thm:Ba} are verifications of these prior results; for
Theorem~\ref{thm:Ba}, rope AV2 discharges \texttt{hsuff} with OpenAI's Lean sufficiency. Our Lean code was
written without consulting theirs.
\item \textbf{Searches.} For Theorems~\ref{thm:C}, \ref{thm:Bb}, \ref{thm:C2} and \ref{thm:D} and
for Conjecture~\ref{conj:M1}, our literature and repository searches (four sweeps, 2026-10-12/13,
over arXiv, GitHub, MathOverflow and bibliographic databases) found no prior statement; for
Theorem~\ref{thm:C} the search log is published with \cite{MR}. The searches for
Theorem~\ref{thm:C2} were made for its circle and two-group forms and were extended in two further
sweeps. On 2026-10-13 a sweep covered Theorem~\ref{thm:C2} in every dimension,
Theorems~\ref{thm:NTG} and~\ref{thm:rank}, Corollary~\ref{cor:radian8} and the prior work credited in
Remark~\ref{rem:definable}. On 2026-10-14 a second sweep covered Theorem~\ref{thm:rigid}; the Clifford
and Vahlen literature by name (Vahlen matrices, Clifford and Lipschitz groups, Ahlfors's Clifford-number
description of M\"obius maps, quaternionic cross-ratios), combined with Ramsey and colouring terms; and
all new items since 2026-10-13. Both used arXiv, OpenAlex, zbMATH, GitHub (including OpenAI's
repository and its forks), the Palomar registry, MathOverflow and Mathematics Stack Exchange, and
neither found a prior statement of these results. The nearest items are P\'alv\"olgyi's appendix
\cite{Palvolgyi} (Lemma~A.5 and Theorems~A.3, A.7 and A.10, which concern certificates and non-Ramsey
criteria and do not give Theorem~\ref{thm:rigid}) and Gibbons and Luo \cite{GibbonsLuo}. One special
case is prior: when the points outside the rigid group include four whose parameters are algebraically
independent over the field generated by three points of the rigid group, the non-Ramsey half of
Theorem~\ref{thm:rigid} follows from P\'alv\"olgyi's Theorem~A.3.
The arXiv search covers titles, abstracts and comments, not full texts; \Kriz's 1992 paper was not
read; the Lean Zulip, Google Scholar and HAL were not searched; and arXiv submissions of 2026-10-14
were not yet indexed. Separately, OpenAI's Proposition~7.3 (first item), which bears on the
conjectures of \S\ref{sec:conj} and the ledger rather than on the theorems above, was overlooked until
2026-10-14, when it was found by reading OpenAI's Lean files. Mirabi \cite{Mirabi} shows that adding to a Ramsey set a point
outside its affine hull keeps it Ramsey, and Moore \cite{Moore} that a pyramid over a Ramsey base is
Ramsey; both add points off the affine hull and do not bear on cocircular sets.
\item \textbf{Name collisions.} ``Two groups'' in this note always means two M\"obius maps, that is, a
union of two M\"obius images of algebraic sets. It is unrelated to the established ``two-orbit''
results in which a soluble group of symmetries acts on the set with at most two orbits (\Kriz{}
\cite{Kriz91}; see Behague \cite{Behague}, who shows that nearly all known Ramsey sets are subsoluble),
and to the ``two-orbit problem'' in \Kriz's sense discussed in an appendix of P\'alv\"olgyi
\cite{Palvolgyi-ncs}.
Gibbons and Luo's ``inversive Ramsey theory'' \cite{GibbonsLuo} colours spheres and asks for
M\"obius copies, not congruent ones; Dam\'asdi, Frankl, Pach and P\'alv\"olgyi \cite{DFPP} study
monochromatic configurations in colourings of a single circle, not Euclidean Ramsey sets in the sense
used here.
Leader--Russell--Walters \cite{LRW12,LRW11} and Graham's survey \cite{Graham94} contain no M\"obius
or cross-ratio remark; \Kriz{} \cite{Kriz91} was read in full and has no M\"obius content;
\cite{Kriz92} could not be obtained. Gibbons and Luo \cite{GibbonsLuo} do not relate the Euclidean
Ramsey property to M\"obius images.
\end{itemize}

\section{Preliminaries}\label{sec:prelim}

We quote OpenAI's classification verbatim from the \TeX{} source at commit \texttt{fd4aeeb2}
(Apache-2.0), as in \cite{MR}; the only edits are cross-reference labels, printed as OpenAI's
rendered numbers. From \cite[\S1.1]{OAI172}:

{\small
\begin{quote}
A singleton is Ramsey. For any other set, let $d\ge1$ be its affine
dimension, and choose a congruent representative
$A=\{a_1,\ldots,a_s\}\subset\R^d$ whose affine span is $\R^d$.
The points are distinct and $s\ge2$. Define
\[
 p_i=\begin{pmatrix}1\\a_i\end{pmatrix},\qquad
 F=\Q\bigl((a_i)_\alpha:1\le i\le s,\ 1\le\alpha\le d\bigr),
 \qquad B=F\otimes_{\Q}F.
\]
The commutative ring $B$ has a multiplication homomorphism
\[
 m_F:B\longrightarrow F,\qquad m_F(x\otimes y)=xy.
\]
An equality in $B$ is stronger than the equality obtained by applying
$m_F$; this map need not be injective. Index the entries of $p_i$ and
the rows and columns of a $(d+1)$-by-$(d+1)$ matrix by $0,1,\ldots,d$.
Thus index $0$ is the constant coordinate.

\begin{oaithm}[Classification]
Let $A=\{a_1,\ldots,a_s\}\subset\R^d$ be a finite set of at least two
points with affine span $\R^d$, and define $p_i,F,B,m_F$ as above.
Then $A$ is Ramsey if and only if there is a matrix
$P\in\operatorname{Mat}_{d+1}(B)$ such that
\begin{align}
 (p_i\otimes1)^T P(1\otimes p_i)&=0
       &&(1\le i\le s),\tag{1}\\
 m_F(P_{\alpha\beta})&=\delta_{\alpha\beta}
       &&(1\le\alpha,\beta\le d).\tag{2}
\end{align}
Here the entries of $p_i\otimes1$ and $1\otimes p_i$ are formed
componentwise. There are no separate conditions on the multiplied
constant row or constant column of $P$.
\end{oaithm}
\end{quote}
}

\noindent\emph{Real form.} Proposition~2.1 of \cite{OAI172} shows that the matrix condition is
equivalent to a tensor condition over $\BB=\R\otimes_\Q\R$. In the array form used in
\cite[\S1]{MR} it reads: there is a $(d+1)\times(d+1)$ array $P$ over $\BB$ with $P^*=P$, where
$P^*_{\alpha\beta}=\tau(P_{\beta\alpha})$ and $\tau$ is the flip, such that
$(p_i\otimes1)^TP(1\otimes p_i)=0$ for all $i$ and the spatial block of $m(P)$ is $I_d$; here
$m:\BB\to\R$ is multiplication. For any $P$ with vanishing evaluations, $\tfrac12(P+P^*)$ keeps them
(apply $\tau$ to an evaluation), and $m(P^*)=m(P)^T$. In OpenAI's Lean proof the array form is
\texttt{RealCriterion}; necessity is \texttt{ramsey\_realCriterion}, descent is
\texttt{realCriterion\_fieldCriterion}, and sufficiency is \texttt{fieldCriterion\_ramsey}.
\textbf{Descent and sufficiency carry no affine-span hypothesis} in Lean; sufficiency needs only
$s\ge2$, $d\ge1$ and injectivity. We use this repeatedly: for a sufficiency argument the spanning
condition never has to be checked.

\noindent\emph{Normalization and scalars.} If $A$ spans and lies on the unit sphere centred at $0$,
a symmetric solution has $m(P)=S:=\mathrm{diag}(-1,I_d)$: the quadratic function $(1,x)m(P)(1,x)^T$
and $\|x\|^2-1$ both vanish on $A$ and differ by an affine function. \emph{Lemma~0}: if $P$ has
vanishing evaluations and the spatial block of $m(P)$ is $cI$ with $c\in\R^\times$ (or
$c\in F^\times$), then $(c^{-1}\otimes1)P$ satisfies the criterion. Ramsey-ness is invariant under
similarities, so every sphere may be taken to be the unit sphere.

\noindent\emph{Circle dictionary.} On the unit circle use the stereographic parameter
$\circmap(t)=((1-t^2)/(1+t^2),\,2t/(1+t^2))$ from the pole $(-1,0)$; a rational rotation moves the
pole off $A$ without changing $F$. Writing $x_i=t_i\otimes1$, $y_i=1\otimes t_i$, a $(2,2)$-form
$f(s,s')=v(s)^TMv(s')$, $v=(1,s,s^2)$, defines $P$ with
$(p_i\otimes1)^TP(1\otimes p_i)=f(x_i,y_i)/((1+x_i^2)(1+y_i^2))$, and $m(f)=\kappa(s-s')^2$ gives
$m(P)=-(\kappa/2)S$. So a circle set is Ramsey iff some $(2,2)$-form over $B$ vanishes at all
$(x_i,y_i)$ and degenerates on the diagonal to $\kappa(s-s')^2$, $\kappa\neq0$.

\section{M\"obius invariance (published)}\label{sec:mobius}

Let $S=\mathrm{diag}(-1,1,\dots,1)$ on $\R^{1+d}$ and $\Ort(1,d)=\{L:L^TSL=S\}$. For $x$ on the
unit sphere $S^{d-1}$ and $L\in\Ort(1,d)$, write $L(1,x)^T=\lambda(1,\varphi_L(x))^T$; the
$0$-th coordinate $\lambda$ is nonzero because $L(1,x)^T$ is a nonzero null vector. For $d\ge2$,
$L\mapsto\varphi_L$ induces the classical isomorphism $\PO(1,d)\cong\Mob(S^{d-1})$,
orientation-reversing maps included \cite{Ratcliffe,Beardon}; see \cite[\S2]{MR} for the exact
locations. We cite this identification; it is not formalized.

\begin{theorem}[Theorem C; published in \cite{MR}]\label{thm:C}
Let $A$ be a finite subset of a sphere in $\R^d$, and let $\varphi$ be a M\"obius transformation of
that sphere. Then $A$ is Ramsey if and only if $\varphi(A)$ is Ramsey. In particular, for points on
a circle, Ramsey-ness depends only on cross-ratios.

\smallskip\noindent\emph{Status:} \st{LEAN} for $s$ distinct points of the unit sphere and
$\varphi=\varphi_L$, $L\in\Ort(1,d)$, with no spanning hypothesis:
\texttt{OAI.EuclideanRamsey.Mobius.}\allowbreak\texttt{ramsey\_iff\_ramsey\_mobius\_full}
(Appendix~\ref{app:lean}). Comparator-checked: rope AV4 (fresh-runner CI run
\texttt{38493016252}, all 12 checks as expected; a fresh-agent reproduction also passed) and the
published package \cite{MR} (three kernels; fresh-runner CI run \texttt{38336666705} on the staging copy, and run
\texttt{38339951762} on the released commit). The package also passed a non-builder fidelity
review and a fresh-agent reproduction (AI; \S\ref{sec:disclosure}), and is registered on the
Palomar registry as PALOMAR-2026-10-14-000008 (version 1, 2026-10-14). The reduction from an arbitrary
sphere to the unit sphere is a similarity \st{HAND} \cite[Remark~1]{MR}, and the identification of
$\Mob(S^{d-1})$ with $\PO(1,d)$ is cited. Scope: \texttt{ANY}.
\end{theorem}

The proof is in \cite[\S2]{MR}. In outline: by Theorem~1.1, the real form and symmetrization, a
Ramsey $A$ has a symmetric certificate $P$ over $\BB$ with $m(P)=S$. With $N=L^{-1}$, the array
$P'=(N\otimes1)^TP(1\otimes N)$ has vanishing evaluations at the image points, because
$Np_i'=\mu_i^{-1}p_i$ with $\mu_i\in\R^\times$, and $m(P')=N^TSN=S$. Descent then gives the field
criterion for the image, whose coordinate field is in general different from $F$; the transport has
to happen over $\BB$, since the entries of $L$ lie in neither field. The project owner read the
hand proof line by line; he describes himself as not an expert in the area, so that read is an
owner's sign-off, not expert review.

\begin{corollary}[{\cite[Cor.~1]{MR}}]\label{cor:mobius}
\begin{enumerate}
\item[(a)] \st{HAND}, scope \texttt{ANY}. Ordered concyclic tuples with equal cross-ratios are both
Ramsey or both not.
\item[(b)] \st{HAND}. Every triangle is Ramsey. This agrees with \cite{FR86} and is immediate from
\cite[Cor.~7.4]{OAI172}; it is not an independent proof.
\item[(c)] \st{LEAN} for the unit sphere (\texttt{ramsey\_mobius\_of\_algebraic},
Comparator-checked in \cite{MR}), \st{HAND} for other spheres and for polygons; scope
\texttt{MOBIUS\_ORBIT\_OF(ALGEBRAIC)}. Every M\"obius image of a finite spherical set with real
algebraic coordinates is Ramsey; in particular every M\"obius image of a regular polygon is.
\item[(d)] \st{LEAN} (\texttt{ramsey\_circ\_mul\_rat}, \texttt{ramsey\_circ\_add\_rat},
Comparator-checked in \cite{MR}). For distinct rationals $t_j$, every real $k\ne0$ and every real
$c$, the sets $\{\circmap(kt_j)\}$ and $\{\circmap(t_j+c)\}$ are Ramsey.
\end{enumerate}
\end{corollary}

For transcendental $k$ the sets in (d) have transcendental coordinate field and are Ramsey for every
number of points, while nine points with algebraically independent parameters are not
\cite[Prop.~7.6]{OAI172}: transcendence alone does not decide the question. Item (c) is the
\emph{one-group} case. Theorem~\ref{thm:C2} below is the two-group case; it does not follow from
Theorem~\ref{thm:C} (Remark~\ref{rem:beyond}).

\section{Algebraic coordinates and circle certificates}\label{sec:algebraic}

\begin{theorem}[Theorem A; prior work \cite{JDJones,Rogerhu,Atimics,Eliott}]\label{thm:A}
A finite spherical set whose coordinates are real algebraic numbers is Ramsey.

\smallskip\noindent\emph{Status:} \st{LEAN}, \textbf{Lean-checked}: rope AV1,
\texttt{ramsey\_of\_algebraic}, with no spanning hypothesis. Stock Comparator under Landlock accepts
all five challenges (including OpenAI's own \texttt{classification} and \texttt{ramsey\_cospherical}
challenges, byte for byte) and refuses all five controls with the expected reasons, locally and on a
fresh CI runner (run \texttt{38493016206}). A non-builder fidelity review found the statement
faithful (its documentation fixes are applied in v1.1, the version in this release), and a
fresh-agent reproduction reproduced it (exit 0, all ten checks as expected; its friction notes are
fixed in the README). Scope:
\texttt{ALGEBRAIC}.
\end{theorem}

\begin{proof}[Proof (as in \cite{JDJones,Rogerhu})]
$F$ is a number field, hence formally unramified over $\Q$, so $B$ contains $e$ with $m(e)=1$ and
$e(1\otimes x-x\otimes1)=0$ for all $x\in F$ (Mathlib
\texttt{Algebra.}\allowbreak\texttt{FormallyUnramified.}\allowbreak\texttt{iff\_\allowbreak{}exists\_\allowbreak{}tensorProduct}). If $h$ is the matrix of the sphere
equation over $F$, then $P=e(h\otimes1)$ satisfies (1), and $m(P)=h$.
\end{proof}

\begin{corollary}[A$'$]\label{cor:Aprime}
\st{HAND}, scope \texttt{ALGEBRAIC} (distances). If all pairwise distances of a finite spherical set
are algebraic, it is Ramsey: Gram--Schmidt over the real closed field $\R\cap\Qbar$ gives a
congruent copy with algebraic coordinates. Hence every non-Ramsey spherical set has a transcendental
distance. JD-Jones's paper text states a closely related remark \cite{JDJones}.
\end{corollary}

\noindent\emph{Graph forms.} A $(1,1)$-form $g(s,s')=\alpha+\beta s+\beta's'+\gamma ss'$ over $B$ is
the graph of a M\"obius correspondence. For three points put
$g_{123}(s,s')=\det\begin{psmallmatrix}ss'&s&s'&1\\ x_ay_a&x_a&y_a&1\end{psmallmatrix}_{a=1,2,3}$;
it vanishes at $(x_a,y_a)$, and $m(g_{123})=\pm(s-s')V(t_1,t_2,t_3)$ with $V$ the Vandermonde.

\begin{theorem}[B(a); prior work \cite{JDJones}]\label{thm:Ba}
Any at most six distinct points on a circle form a Ramsey set. With P\'alv\"olgyi's heptagon,
seven is the minimum size of a non-Ramsey set of concyclic points.

\smallskip\noindent\emph{Status:} \st{LEAN}, Comparator-checked: rope AV2,
\texttt{circle\_ramsey\_of\_le\_six} (fresh-runner CI run \texttt{38493016157}, all ten checks as
expected); a fresh-agent reproduction passed (exit 0, all ten checks as expected), and the rope's
non-builder fidelity review is still open. Also \st{CERT}
(\texttt{certificates/cert\_mobius.py}). Scope: \texttt{ANY}. The certificate $f=g_{123}g_{456}$ is
JD-Jones's two-triple certificate.
\end{theorem}

\begin{theorem}[B(b)]\label{thm:Bb}
Let $A$ be a finite set of points on a circle in $\R^2$, all but at most three of which have real
algebraic coordinates. Then $A$ is Ramsey. (The centre and radius are not assumed algebraic.)

\smallskip\noindent\emph{Status:} \st{LEAN}, \textbf{Lean-checked}: rope AV5,
\texttt{ramsey\_alg\_plus\_three} (fresh-runner CI run \texttt{38493016164}, all ten checks as
expected). A non-builder fidelity review found the statement faithful (documentation fixes
applied), and a fresh-agent reproduction reproduced it (exit 0, all ten checks as expected; its
friction notes are fixed in the documentation). Scope: \texttt{ALGEBRAIC} $\cup$
$3\cdot$\texttt{ANY}. The Lean statement is the theorem as stated: an injective family of points of
some circle in $\R^2$ (centre and radius arbitrary), all but at most three of which have algebraic
coordinates. No prior statement found in our searches. It settles the class that P\'alv\"olgyi's
A.10 \cite{Palvolgyi} leaves undecided (A.10 shows only that derivation certificates cannot
decide it). It is not of the form of \cite{Mirabi,Moore}, where the added points lie off the affine
hull: here every point lies on the circle.
\end{theorem}

\begin{proof}[Idea]
Take $f=(s-s')\,g_{123}\cdot e_K$, where $g_{123}$ passes through the three arbitrary points and
$e_K\in K\otimes K$ is the separability idempotent of the number field $K$ generated by the
algebraic parameters, so that $e_K(x_i-y_i)=0$ at every algebraic point. For $s\le6$ use
Theorem~\ref{thm:Ba}; for $s\ge7$ at least four points are algebraic, which makes the centre and
radius algebraic.
\end{proof}

The circle case of the two-group theorem is Theorem~\ref{thm:C2} with $d=2$; Theorems~\ref{thm:Ba} and \ref{thm:Bb} are special cases of it in which one group is a triple
(every triple is a M\"obius image of $\{0,1,\infty\}$), although their Lean proofs are independent of
it.

\section{Two M\"obius groups in every dimension}\label{sec:twogroups}

\begin{theorem}[Theorem C2, the two-group theorem]\label{thm:C2}
Let $d\ge2$, let $A_1',A_2'\subset S^{d-1}\subset\R^d$ be finite sets whose coordinates are real
algebraic numbers, and let $\nu_1,\nu_2$ be M\"obius transformations of $S^{d-1}$
(orientation-reversing ones allowed). Then $\nu_1(A_1')\cup\nu_2(A_2')$ is Ramsey. There is no bound
on $|A_1'|$ or $|A_2'|$.

\smallskip\noindent\emph{Status:}
\begin{itemize}
\item $d\le6$ ($S^1$ to $S^5$): \st{CERT} plus written glue. The identity scripts were re-run
independently with identical output; adversarial reviews (AI) found the glue correct, with wording fixes
incorporated here: for $d\le5$ the complex and quaternionic forms (\S\ref{sec:lowdim}), for $d=6$
the Clifford form of \S\S\ref{sec:clifford}--\ref{sec:seeds}.
\item every $d\ge2$: \st{HAND}, reviewed (AI) as correct; the identity families are
machine-checked symbolically for small $d$ and on exact random instances through $d=8$
(\S\ref{sec:lowdim}).
\item Coverage of every M\"obius map rests on two cited textbook facts: the decomposition of a
M\"obius map into an inversion and a similarity \cite{Beardon}, and the Cartan--Dieudonn\'e
theorem \cite{CartanDieudonne}.
\end{itemize}
Scope: \texttt{MOBIUS\_ORBIT\_OF(ALGEBRAIC)}$\times2$. No prior statement found in our searches (see
\S\ref{sec:credit} for their coverage). ``Two groups'' means two M\"obius maps;
the theorem is unrelated to the soluble ``two-orbit'' results of \cite{Kriz91,Behague}
(\S\ref{sec:credit}).
\end{theorem}

\begin{remark}[What two groups add]\label{rem:beyond}
\st{HAND} Theorem~\ref{thm:C2} is not a consequence of Theorem~\ref{thm:C} and
Theorem~\ref{thm:A}. On a circle, any four points of a M\"obius image of an algebraic set have an
algebraic cross-ratio. In stereographic parameters, the union of the rational set $\{0,1,\infty\}$
(group~1, $\nu_1=\mathrm{id}$) and $\{c\}$ with $c$ transcendental (group~2, $\nu_2(t)=t+c$,
$A_2'=\{0\}$) has the cross-ratio $c$ (up to the usual six substitutions, all transcendental), so it is not a
M\"obius image of an algebraic set. So the two-group class is strictly larger than the one-group class
of Corollary~\ref{cor:mobius}(c); it contains sets with arbitrarily many points and transcendental
cross-ratios, whereas nine points with independent parameters are not Ramsey
\cite[Prop.~7.6]{OAI172}.
\end{remark}

\begin{remark}[Relation to prior work]\label{rem:six}
\st{HAND} M\"obius maps act transitively on ordered triples of distinct points of $S^{d-1}$, so
any set of at most six points on a sphere is a union of two M\"obius images of rational triples, and
Theorem~\ref{thm:C2} makes it Ramsey. This corollary is weaker than JD-Jones's $m+4$ theorem
\cite{JDJones} together with OpenAI's sufficiency, and is not claimed.
\end{remark}

\begin{remark}[Three free points in every dimension]\label{rem:C4}
At the level of Theorem~\ref{thm:C2} (\st{CERT}+glue for $d\le6$, \st{HAND} for all $d$). Let
$A\subset S^{d-1}$, $d\ge2$, consist of finitely many points with real algebraic coordinates
together with at most three further points, arbitrary. The algebraic points are one group
($\nu_1=\mathrm{id}$), and the further points are the image of a rational set of the same size
under a M\"obius map (M\"obius maps act transitively on ordered triples of distinct points), so $A$
is Ramsey. This is the first half of Conjecture C4 (\S\ref{sec:conj}) on every unit sphere, and the
spherical analogue of Theorem~\ref{thm:Bb}. Four free points can fail: on $S^1$ by
Proposition~\ref{prop:PD}, on $S^2$ and $S^3$ by the cells $(6,4)$ and $(11,4)$ of
Theorem~\ref{thm:NR}. The two-line derivation has had no separate review.
\end{remark}

\subsection{Clifford conventions and coordinates}\label{sec:clifford}

Put $n=d-1$ (the dimension of the stereographic chart) and let $\Cl_n:=\Cl_{0,n-1}$ be the real
associative algebra generated by $e_1,\dots,e_{n-1}$ with $e_ie_j+e_je_i=-2\delta_{ij}$; it has
dimension $2^{n-1}$ and a basis of blades with structure constants $\pm1$. ($\Cl_1=\R$, $\Cl_2\cong\C$,
$\Cl_3\cong\HH$.) For any commutative $\Q$-algebra $R$ let $\Cl_n(R)$ be the $R$-algebra with the
same structure constants; every identity below is a polynomial identity with integer coefficients
in blade coordinates, so it holds in $\Cl_n(R)$ for every such $R$. \emph{Paravectors} are the
elements $x=x_0+x_1e_1+\dots+x_{n-1}e_{n-1}$; we identify them with $\R^n$. We use the reversion
$a\mapsto a^*$ (the anti-automorphism fixing each $e_i$), the grade involution $a\mapsto\hat a$
(the automorphism $e_i\mapsto-e_i$) and Clifford conjugation $\bar a=\hat a^{\,*}$. For a paravector
$x$: $x^*=x$, $\hat x=\bar x$, and $x\bar x=\bar xx=|x|^2$. The scalar part $\Rea$ satisfies
$\Rea(ab)=\Rea(ba)$ and $\Rea(a\bar b)=\langle a,b\rangle$, the Euclidean inner product in blade
coordinates.

\noindent\emph{Coordinates.} For a paravector $x$ put
\[
 a(x)=\frac{(2x,\ |x|^2-1)}{1+|x|^2}\in S^{d-1},\qquad
 b(x)=(1,x_0,\dots,x_{n-1},|x|^2),\qquad q(x)=(1+|x|^2,\ 2x,\ |x|^2-1).
\]
This is the inverse stereographic projection from the pole $\omega=(0,\dots,0,1)$. Then $q=Cb$ for an
invertible rational matrix $C$, the OAI vector of the point $a(x)$ is $p(x)=(1,a(x))=q(x)/(1+|x|^2)$,
and $q(x)^TSq(y)=-2|x-y|^2$. A bilinear form $f(x,y)=b(x)^TMb(y)$ corresponds to the array
$P=C^{-T}MC^{-1}$, with $p(x)^TPp(y)=f(x,y)/((1+|x|^2)(1+|y|^2))$; since the entries of $b$ are
linearly independent polynomials, $M$ is determined by $f$.

\begin{lemma}[Pairing]\label{lem:pairing}
Let $R$ be a commutative $\Q$-algebra and, for $k=1,2$, let
$h_k(x,y)=y\gamma_kx+y\delta_k+\alpha_kx+\beta_k$ with $\alpha_k,\beta_k,\gamma_k,\delta_k\in\Cl_n(R)$.
Then $\Rea\bigl(h_1(x,y)\,\overline{h_2(x,y)}\bigr)=b(x)^TM\,b(y)$ for paravectors $x,y$ with
coordinates in $R$, where $M\in\Mat_{d+1}(R)$ depends bilinearly on the coefficients of $h_1,h_2$ and
not on $x,y$.
\end{lemma}

\begin{proof}
\st{HAND}; machine-checked for $n=2,\dots,7$. Expand
$h_1\bar h_2$ with $\bar h_2=\bar x\bar\gamma_2\bar y+\bar\delta_2\bar y+\bar x\bar\alpha_2+\bar\beta_2$.
In each term $x$ is rightmost in its factor of $h_1$ and $\bar x$ leftmost in its factor of $\bar
h_2$, so $x$ meets $\bar x$ only as $x\bar x=|x|^2$, which is central. The variable $y$ occurs at most
once on each side, and $\Rea(y\,u\,\bar y)=|y|^2\Rea(u)$ by cyclicity. So every term is $\Rea$ of a
product that is affine in $x$ (or contains $|x|^2$) and affine in $y$ (or contains $|y|^2$), and
lies in $\mathrm{span}\{1,x_i,|x|^2\}\otimes\mathrm{span}\{1,y_j,|y|^2\}$.
\end{proof}

\subsection{Lipschitz elements and the Vahlen normal form}

Vahlen matrices \cite{Vahlen,Ahlfors85,EGM,Waterman} describe the M\"obius maps of $\R^n$ by $2\times2$ matrices over a
Clifford algebra. We do not use their characterization, only the explicit normal form below, whose
coverage we prove directly (Lemma~\ref{lem:coverage}).

\begin{lemma}[Lipschitz]\label{lem:lipschitz}
\st{HAND}; machine-checked for $n=1,\dots,7$. For paravectors $p,s$: $p\bar p=\bar pp=|p|^2$ and
$psp=2\langle p,\bar s\rangle p-|p|^2\bar s$, a paravector. Hence for $w=p_1\cdots p_\ell$ a product of
paravectors, $w\bar w=\bar ww=\hat ww^*=w^*\hat w=|w|^2:=\prod|p_i|^2$, a scalar, and $s\mapsto
wsw^*$ maps paravectors to paravectors. If $\ell$ is even and $|w|\neq0$, then
$R_w(s):=ws\hat w^{-1}=wsw^*/|w|^2$ is the composite of the $\ell$ reflections in the hyperplanes
$p_1^\perp,\bar p_2^\perp,p_3^\perp,\dots$, an element of $\SO(n)$.
\end{lemma}

\begin{proof}
Polarizing $x\bar x=|x|^2$ gives $s\bar p+p\bar s=2\langle s,p\rangle$; replacing $s$ by $\bar s$
and multiplying by $p$ gives the formula for $psp$, which equals $-|p|^2\rho_p(\bar s)$ with $\rho_p$
the reflection in $p^\perp$. Since $w^*=p_\ell\cdots p_1$ and $\hat w=\bar p_1\cdots\bar p_\ell$, the norm
identities telescope, and $wsw^*$ is an iterated $psp$. Pairs of conjugations cancel, because
conjugation $\kappa(s)=\bar s$ is orthogonal and $\kappa\rho_p\kappa=\rho_{\bar p}$.
\end{proof}

For $w$ as above with $\ell=2\lfloor n/2\rfloor$ factors, $\lambda\in\R^\times$ and paravectors $p,q$,
put
\[
 N=\begin{pmatrix}q\hat w&\lambda w+q\hat wp\\ \hat w&\hat wp\end{pmatrix},\qquad
 \nu(t)=(N_{11}t+N_{12})(N_{21}t+N_{22})^{-1}=q+\lambda\,w(t+p)^{-1}\hat w^{-1}.
\]
For a $2\times2$ matrix $G$ over $\Cl_n(R)$ put
$\adj G=\begin{psmallmatrix}G_{22}^*&-G_{12}^*\\-G_{21}^*&G_{11}^*\end{psmallmatrix}$.

\begin{lemma}[Vahlen identities]\label{lem:vahlen}
\st{HAND}; machine-checked symbolically for $n=2,3,4,5$. With $u:=N_{21}t+N_{22}=\hat w(t+p)$:
\begin{enumerate}
\item[(V1)] $u\bar u=\bar uu=|t+p|^2|w|^2=:\mathrm n(t)$, a scalar;
\item[(V2)] $\tilde X:=(N_{11}t+N_{12})\bar u=q\,\mathrm n(t)+\lambda\,w\,\overline{(t+p)}\,w^*$ is a
paravector, so $\nu(t)=\tilde X/\mathrm n(t)$ lies in the chart;
\item[(V3)] $\adj(N)N=N\adj(N)=\Delta I$ with $\Delta=-\lambda|w|^2\in\R^\times$.
\end{enumerate}
\end{lemma}

\begin{proof}
(V1) and (V2) follow from Lemma~\ref{lem:lipschitz}, since $\bar{\hat w}=w^*$ and $N_{11}t+N_{12}=qu+\lambda
w$. For (V3) expand the eight entries; each collapses, by $\hat w^*=\bar w$ and the norm identities of
Lemma~\ref{lem:lipschitz}, to $-\lambda w^*\hat w$, $0$, $0$, $-\lambda\bar ww$ (for $\adj(N)N$) and
to $-\lambda w\bar w$, $0$, $0$, $-\lambda\hat ww^*$ (for $N\adj(N)$). For example
$(\adj(N)N)_{12}=\lambda(p\,\bar ww-w^*\hat w\,p)=0$. Nothing uses $n$.
\end{proof}

\subsection{Graph spinors and vanishing}

Let $R=\BB$. For a normal-form matrix $N$ with real entries write $N^{(1)}$ and $N^{(2)}$ for its
images under $r\mapsto r\otimes1$ and $r\mapsto1\otimes r$, applied to blade coordinates. Put
\[
 \mu=N^{(2)}\adj(N^{(1)}),\qquad
 h(x,y)=y\,\mu_{21}\,x+y\,\mu_{22}-\mu_{11}\,x-\mu_{12},
\]
the \emph{graph form} of the M\"obius map $\nu$ of $N$.

\begin{lemma}[Vanishing]\label{lem:vanish}
\st{HAND}; reviewed (AI) as correct. Let $t$ be a paravector with rational coordinates and
$t\neq-p$. Then $h\bigl(\nu(t)\otimes1,\,1\otimes\nu(t)\bigr)=0$ in $\Cl_n(\BB)$.
\end{lemma}

\begin{proof}
Because $t$ is rational, $\nu(t)\otimes1=\tilde X/\mathrm n$ and $1\otimes\nu(t)=\tilde Y/\mathrm n'$,
where $\tilde X,\mathrm n$ are computed from $N^{(1)}$ and $\tilde Y,\mathrm n'$ from $N^{(2)}$ at the same
$t$; $\mathrm n=\mathrm n(t)\otimes1$ and $\mathrm n'=1\otimes\mathrm n(t)$ are pure tensors of nonzero
reals, hence central units. Expanding,
\[
 \mathrm n\,\mathrm n'\,h(\tilde X/\mathrm n,\tilde Y/\mathrm n')
 =\begin{pmatrix}-\mathrm n'&\tilde Y\end{pmatrix}\mu\begin{pmatrix}\tilde X\\ \mathrm n\end{pmatrix}.
\]
With $u=N^{(1)}_{21}t+N^{(1)}_{22}$, the definitions and (V1) give $\begin{psmallmatrix}\tilde X\\ \mathrm n\end{psmallmatrix}
=N^{(1)}\begin{psmallmatrix}t\\1\end{psmallmatrix}\bar u$, so by (V3)
$\mu\begin{psmallmatrix}\tilde X\\ \mathrm n\end{psmallmatrix}
=\Delta^{(1)}N^{(2)}\begin{psmallmatrix}t\\1\end{psmallmatrix}\bar u$, with $\Delta^{(1)}=\Delta\otimes1$
central. Finally, with $u'=N^{(2)}_{21}t+N^{(2)}_{22}$,
\[
 \begin{pmatrix}-\mathrm n'&\tilde Y\end{pmatrix}N^{(2)}\begin{pmatrix}t\\1\end{pmatrix}
 =-\mathrm n'(N^{(2)}_{11}t+N^{(2)}_{12})+(N^{(2)}_{11}t+N^{(2)}_{12})\,\bar u'u'=0,
\]
by the other order of (V1). All products are in the associative algebra $\Mat_2(\Cl_n(\BB))$.
\end{proof}

Applying $m$ coefficientwise (a ring map that commutes with the involutions) gives
$m(\mu)=N\adj(N)=\Delta I$, so $m(h)(x,y)=\Delta\,(y-x)$.

\subsection{Coverage}

\begin{lemma}[Every M\"obius map, up to rational changes of seeds]\label{lem:coverage}
\st{HAND}, reviewed (AI) as correct, modulo the two cited facts. Let $\nu$ be a M\"obius
transformation of $\R^n\cup\{\infty\}$, of either orientation, and $A'\subset\R^n\cup\{\infty\}$ finite,
with rational (respectively algebraic) points, such that $\infty\notin\nu(A')$. Then $\nu(A')=\nu'(A'')$
for a finite $A''\subset\R^n$ with rational (respectively algebraic) points and a map $\nu'$ in the
normal form above; and $\nu'(t)\neq\infty$ for $t\in A''$.
\end{lemma}

\begin{proof}
If $\nu(\infty)=\infty$ or $\infty\in A'$, replace $(\nu,A')$ by $(\nu\circ\rho,\rho(A'))$, where $\rho$
is the inversion in a sphere with rational centre $r\notin A'$, $r\ne\nu^{-1}(\infty)$, and rational
radius; only finitely many $r$ are excluded, $\rho$ is an involution defined over $\Q$, $\rho(A')\subset\R^n$,
and $\nu\circ\rho(\infty)=\nu(r)\ne\infty$. So assume $\nu(\infty)\neq\infty$ and let
$c=\nu^{-1}(\infty)$. Then $\psi=\nu\circ T_c\circ\sigma$, with $\sigma(x)=x/|x|^2$ and $T_c$ the
translation by $c$, fixes $\infty$ and is therefore a similarity $x\mapsto Lx+q$ with
$L\in\CO(n)=\R_{>0}\cdot\Ort(n)$ \cite[\S3.5]{Beardon}; hence $\nu(x)=q+L\sigma(x-c)$. The normal form
is $q+\lambda R_w\kappa\sigma(x+p)$, since $(x+p)^{-1}=\kappa\sigma(x+p)$; take $p=-c$. By
Cartan--Dieudonn\'e every element of $\SO(n)$ is a product of an even number $\le n$ of reflections,
hence (padding with $\rho\rho=1$) of exactly $2\lfloor n/2\rfloor$, so $\{R_w\}=\SO(n)$ by
Lemma~\ref{lem:lipschitz}. Conjugation $\kappa$ negates $n-1$ coordinates. If $n$ is odd,
$\kappa\in\SO(n)$ and $-I\notin\SO(n)$, so $\{\lambda R_w\kappa\}=\R^\times\SO(n)=\CO(n)$ and every
$\nu$ is covered. If $n$ is even, $\det\kappa=-1$ and $-I\in\SO(n)$, so the family consists of exactly the
$L\in\CO(n)$ with $\det L<0$; if $\det L>0$, replace $(\nu,A')$ by $(\nu\circ\varepsilon,\varepsilon(A'))$
for a coordinate reflection $\varepsilon$, whose linear part is $L\varepsilon$. ($n=1$ is the degenerate case $w=1$,
$\CO(1)=\R^\times$.)
\end{proof}

\subsection{Proof of Theorem~\ref{thm:C2}}\label{sec:proofC2}

\st{HAND} for every $d$; \st{CERT} for $d\le6$ in the sense of \S\ref{sec:lowdim}. Let
$A=\nu_1(A_1')\cup\nu_2(A_2')$.

\emph{Step 1 (configuration).} Rational rotations are dense in $\SO(d)$ (Cayley transform), so
choose one, $\rho_0$, with $\omega\notin\rho_0A$. Since $\rho_0A=(\rho_0\nu_1)(A_1')\cup(\rho_0\nu_2)(A_2')$ is
congruent to $A$, we may assume $\omega\notin A$. Pass to the chart by stereographic projection from
$\omega$, which is defined over $\Q$: rational points of $S^{d-1}$ other than $\omega$ correspond to
rational chart points, a seed at $\omega$ to $\infty$, and the $\nu_k$ to M\"obius maps of $\R^n\cup\{\infty\}$. By
Lemma~\ref{lem:coverage}, each $\nu_k$ may be taken in normal form, with finite rational seed sets
$T_k\subset\R^n$, and $\nu_k(t)\neq\infty$ for $t\in T_k$. This step keeps rational seeds
rational and algebraic seeds algebraic.

\emph{Step 2 (certificate).} Let $h_k$ be the graph form of $\nu_k$, put
$f=\Rea(h_1\bar h_2)=b(x)^TMb(y)$ by Lemma~\ref{lem:pairing} with $R=\BB$, and
$P=C^{-T}MC^{-1}\in\Mat_{d+1}(\BB)$.

\emph{Step 3 (evaluations).} A point of $A$ is $a(\nu_k(t))$ with $t\in T_k$; its OAI vector is
$p=(1,a(\nu_k(t)))$, in OpenAI's coordinates on the unit sphere. Its evaluation is
$(p\otimes1)^TP(1\otimes p)=f(x,y)/((1+|x|^2)\otimes(1+|y|^2))$ with $x=\nu_k(t)\otimes1$,
$y=1\otimes\nu_k(t)$, where the denominator is a pure tensor of nonzero reals. By
Lemma~\ref{lem:vanish}, $h_k(x,y)=0$, so $f(x,y)=0$ (for $k=2$ use $\overline{0}=0$).

\emph{Step 4 (normalization).} Since $m$ is a ring map, $m(f)=\Rea\bigl(\Delta_1(y-x)\,\Delta_2\overline{(y-x)}\bigr)
=\Delta_1\Delta_2|x-y|^2=-\tfrac12\Delta_1\Delta_2\,q(x)^TSq(y)$, so $m(P)=-\tfrac12\Delta_1\Delta_2\,S$
with $\Delta_k=-\lambda_k|w_k|^2\in\R^\times$. Apply Lemma~0 with $c=-\tfrac12\Delta_1\Delta_2$ and then
replace the result by its symmetrization $\tfrac12(P+P^*)$: this is OpenAI's real criterion for $A$. OpenAI's Lean descent and
sufficiency (\S\ref{sec:prelim}) need no spanning hypothesis, so $A$ is Ramsey. \qed

\subsection{Algebraic seeds}\label{sec:seeds}

\st{CERT}+glue, reviewed (AI) as correct, and checked exactly on a $\Q(\sqrt2,\sqrt3)$ example on the
circle (\texttt{certificates/g4\_test.py}). Let $K\subset\R$ be a number field containing every seed
coordinate of both groups and $e_K\in K\otimes_\Q K\subset\BB$ its separability idempotent, so that
$m(e_K)=1$ and $e_K(r\otimes1-1\otimes r)=0$ for $r\in K$. With $t_L,t_R$ independent paravector
variables, $\Phi(t_L,t_R):=\mathrm n\,\mathrm n'\,h(\nu^{(1)}(t_L),\nu^{(2)}(t_R))$ is a polynomial in the
coordinates of $t_L,t_R$ with coefficients in $\Cl_n(\BB)$, and Lemma~\ref{lem:vanish} (whose proof is
an identity in a symbolic $t$) says that it vanishes on $t_L=t_R$. So $\Phi$ lies in the ideal
generated by the $n$ coordinates of $t_L-t_R$, and $e_K\Phi(r\otimes1,1\otimes r)=0$ for every algebraic
seed $r$; the denominators stay pure tensors. Hence $e_KP$ has vanishing evaluations at every point of
both groups, and $m(e_KP)=m(P)$. Step~4 applies unchanged.

\subsection{Exact certificates for $d\le6$, and cross-checks}\label{sec:lowdim}

\noindent\emph{Circle and $S^2$ ($d=2,3$), complex form.} On $S^2$ the chart is $\C$ and the graph
forms are $\C$-bilinear, $g_k=\bigl(\adj N_k^{(1)}(z,1)^T\bigr)\wedge\bigl(\adj N_k^{(2)}(w,1)^T\bigr)$ for
$\nu_k\leftrightarrow N_k\in\PGL_2(\C)$ (orientation-reversing maps are $\nu_k$ composed with
conjugation). The Hermitian product $H=g_1\,\overline{g_2}$ gives two real
certificates $\Rea H$ and $\operatorname{Im}H$, with $m(H)=W_1\overline{W_2}|z-w|^2$, $W_k=\det N_k$;
since the scalar $W_1\overline{W_2}$ is complex, the na\"ive choice $\Rea H$ alone can degenerate,
and one of $\Rea H$, $\operatorname{Im}H$ (or $\overline{W_1}W_2\cdot H$) always works.
\texttt{certificates/cert\_s2.py} checks the identities symbolically in a seed point, for generic complex
parameters in each tensor factor, including an explicit $6+6$-point instance, and checks a third group
as a control (nonzero). The circle case follows by specializing to real maps (imaginary parameters
$0$, $v=0$).

\noindent\emph{$S^3$, $S^4$ ($d=4,5$), quaternionic form.} \texttt{certificates/cert\_s2\_quat.py} checks
Lemma~\ref{lem:pairing} over $\HH$, the diagonal, and the graph form of a general quaternionic
M\"obius map $\nu(t)=(t+\beta)(\gamma t+\delta)^{-1}$ (with a general leading coefficient), whose
diagonal is a positive multiple of $y-x$. Coverage there was argued through the surjectivity of
$\PSL_2(\HH)$ onto the orientation-preserving M\"obius group of $S^4$ \cite{Wilker} (reviewed:
correct, with that fact cited); $S^3$ is a great $3$-sphere of $S^4$. The Clifford route of this
section covers $d=4,5$ as well, with coverage by Lemma~\ref{lem:coverage} instead.

\noindent\emph{$S^5$ ($d=6$), Clifford form.} \texttt{certificates/cert\_clifford.py} checks in
$\Cl_{0,4}$: the algebra axioms; Lemma~\ref{lem:pairing} with 128 free coefficient symbols; the
diagonal; Lemma~\ref{lem:lipschitz}; Lemma~\ref{lem:vahlen} symbolically (36 symbols); and an
end-to-end test at exact rationals with independent parameters on the two sides and five seeds per
group, in OpenAI's coordinates, with a third group as a control (54--63 minutes).
\texttt{cert\_clifford\_seed.py} checks the input to \S\ref{sec:seeds}.

\noindent\emph{The naive route fails.} \st{CERT} $\Cl_{0,4}$ has zero divisors,
$(1+e_{1234})(1-e_{1234})=0$, and the generic-entry recipe of the quaternionic case breaks there: the
invertible non-Lipschitz element $2+e_{123}$ sends the paravector $1$ off the chart under $t\mapsto
t(2+e_{123})^{-1}$, and used as $w$ it breaks (V1)--(V3) (\texttt{cert\_clifford\_obs.py}). This is
why the normal form restricts $w$ to the Lipschitz group.

\noindent\emph{Every $d$.} Only coverage depends on $n$ (the number of factors of $w$ and the parity
of $\det\kappa=(-1)^{n-1}$). Lemma~\ref{lem:pairing} is machine-checked for $n=2,\dots,7$,
Lemma~\ref{lem:lipschitz} for $n=1,\dots,7$, and Lemma~\ref{lem:vahlen} symbolically for $n=2,\dots,5$.
An independent script with its own blade-sign table, exact fractions and inverses by linear solving
(\texttt{certificates/cert\_clifford\_crosscheck.py}) checks (V1)--(V3), the vanishing of
Lemma~\ref{lem:vanish} and the pairing identity on exact random instances for $n=2,\dots,7$
($d=3,\dots,8$), and a floating-point test recovers every M\"obius map with $\nu(\infty)\ne\infty$ in
the normal-form family for $n=2,\dots,7$, both orientations for odd $n$ and exactly one for even $n$.

\begin{remark}[Why two]\label{rem:whytwo}
\st{SKETCH} $f$ is bilinear and symmetric in $(\mu_1,\mu_2)$, so it is the polar form of the quadratic
map $\mu\mapsto f(\mu,\mu)$, and for one Vahlen matrix $f(\mu,\mu)$ is, up to scale, the Lorentz
matrix of $\mu$ in $\Ort(1,d)$ (tested on three rational instances on $S^5$,
\texttt{cert\_clifford\_lorentz.py}). A polar form has exactly two slots. The words ``spinor'' and
``Lorentz'' are heuristic here: $\mu$ is not a group element, and the matrix is Lorentz only up to
scale. This explains why the construction stops at two groups; it is \emph{not} evidence that three
groups fail. The evidence for that is separate: three M\"obius-algebraic groups fail the second-order
test of \S\ref{sec:jet} in every tested case on $S^2$, $S^3$, $S^4$ (\st{PROBE}).
\end{remark}

\noindent\emph{Open.} Three groups (Conjecture~\ref{conj:R1} for the circle). A single Lean statement
uniform in $d$, over Mathlib's \texttt{CliffordAlgebra}, would need the even-$n$ orientation step
inside the family.

\section{Radian progressions and transcendence}\label{sec:radian}

\begin{theorem}[Theorem D]\label{thm:D}
Let $\theta\in\R$ be such that the nine points $E_\theta=\{(\cos k\theta,\sin k\theta):0\le k\le8\}$
are distinct. Then $E_\theta$ is Ramsey if and only if $e^{i\theta}$ is algebraic (equivalently,
$\tan(\theta/2)$ is algebraic; equivalently, $\cos\theta$ and $\sin\theta$ are).

\smallskip\noindent\emph{Status:} \st{LEAN}, \textbf{Lean-checked}: rope AV3 (fresh-runner CI run
\texttt{38493016161}, all ten checks as expected), theorems \texttt{radian\_\allowbreak{}configuration} (if
$\tan(\theta/2)$ is transcendental, the nine points are distinct, on the unit circle, and not Ramsey),
\texttt{radian\_\allowbreak{}not\_\allowbreak{}ramsey}, \texttt{radian\_\allowbreak{}ramsey\_\allowbreak{}iff} and \texttt{radian\_\allowbreak{}ramsey\_\allowbreak{}iff\_\allowbreak{}tan};
a non-builder fidelity review found the statements faithful (documentation fixes applied), and a
fresh-agent reproduction reproduced them (exit 0, all ten checks as expected). The equivalence with $e^{i\theta}$ is a one-line
\st{HAND} step. Scope: \texttt{TRANSCENDENTAL\_\allowbreak{}PARAM}$(\tan\theta/2)$ / \texttt{ALGEBRAIC}. No prior
statement found in our searches; \cite[Prop.~7.6]{OAI172} is the case of nine independent
parameters.
\end{theorem}

\begin{proof}[Proof sketch]
($\Leftarrow$) Theorem~\ref{thm:A}. ($\Rightarrow$) If $t=\tan(\theta/2)$ is transcendental then
$F=\Q(t)$ and $B$ embeds in $\Q(T,U)$. The $9\times9$ evaluation matrix of condition (1) has nonzero
determinant in $\Q(T,U)$ (nonzero at the rational points $(1/3,2/7)$ and $(-5/11,3/13)$, \texttt{certificates/radian\_det.py}; the Lean proof certifies the denominator-cleared determinant at $(T,U)=(2,3)$ by an integer triangular factorization), so (1)
forces $P=0$, contradicting (2). This is the argument of \cite[Prop.~7.6]{OAI172}, and the Lean proof
mirrors OpenAI's \texttt{EuclideanRamseyNine}.
\end{proof}

\begin{corollary}\label{cor:radian}
\begin{enumerate}
\item \st{HAND} (Theorem~\ref{thm:D} plus Lindemann), scope \texttt{TRANSCENDENTAL\_}\allowbreak\texttt{PARAM}$(\tan\tfrac12)$.
The nine points at $0,1,\dots,8$ radians are not Ramsey, since $e^{i}$ is transcendental.
(Lindemann--Weierstrass is not in Mathlib \texttt{d13f23b7}, so this instance is not in Lean.)
\item \st{HAND} (Theorem~\ref{thm:D}; distinctness below), scope as in Theorem~\ref{thm:D}.
The nine points at $0,e,\dots,8e$ radians are Ramsey iff $e^{ie}$ is algebraic. (The points are
distinct because $ke\notin2\pi\mathbb Z$ for $1\le k\le8$: since $ke\le8e<4\cdot2\pi$, this only
excludes the finitely many values $e/\pi=2j/k$ with $1\le j\le3$, $1\le k\le8$, which a numerical
evaluation of $e/\pi$ does.) We know no
proof that $e^{ie}$ is transcendental and found none in the literature; that the question is open is
our inference, from the fact that Schanuel's conjecture implies it \cite[\S3.1]{Waldschmidt}: if
$e^{ie}$ were algebraic, then $\Q(1,ie,e,e^{ie})$ would have transcendence degree at most $1<2$
(literature re-checked on 2026-10-14). So explicit Euclidean Ramsey questions can be as hard as open
transcendence problems (Conjecture~\ref{conj:T1}).
\item \st{HAND} (Theorem~\ref{thm:D} plus Gelfond--Schneider \cite{Baker}; not reviewed), scope
\texttt{ALGEBRAIC} in $\beta$. If $\theta=\pi\beta$ with $\beta$ real algebraic and the nine points
distinct, they are Ramsey iff $\beta\in\Q$: then $e^{i\theta}$ is a root of unity, while for irrational
algebraic $\beta$ the number $e^{i\pi\beta}=(-1)^\beta$ is transcendental by Gelfond--Schneider.
\end{enumerate}
\end{corollary}

Eight consecutive multiples are Ramsey for every $\theta$ (Corollary~\ref{cor:radian8} below), so nine is
the exact threshold for radian progressions with transcendental $e^{i\theta}$.

\section{Beyond two groups on the circle}\label{sec:beyond}

\begin{remark}[Quadric independence; prior work \cite{OAI172}]\label{rem:quadric}
OpenAI's Proposition~7.3 \cite{OAI172} (\st{LEAN}, theirs: \texttt{spherical\_\allowbreak{}quadraticIndependent\_\allowbreak{}ramsey\_\allowbreak{}all},
with spanning hypothesis, and \texttt{spherical\_\allowbreak{}quadraticIndependent\_\allowbreak{}ramsey\_\allowbreak{}of\_\allowbreak{}dimension}, without
one) says that a finite spherical set is Ramsey if the vectors $p_ip_i^T$ are linearly independent over
$F$, that is, if the points impose independent conditions on quadratic polynomials. Quadratic
polynomials restrict to a space of dimension $D_d=(d+1)(d+2)/2-1$ on $S^{d-1}$. Corollary
(\st{LEAN}, OpenAI's, plus a \st{HAND} genericity step reviewed once by an AI agent, plus exact ranks;
scope \texttt{GENERIC\_OF}): for every $d\ge2$, any $n\le D_d$ points of $S^{d-1}$ whose stereographic
coordinates are algebraically independent over $\Q$ are Ramsey. Indeed $D_d$ such points affinely span
$\R^d$ and have independent rows: quadratic polynomials restrict to linearly independent functions on
$S^{d-1}$, so the $D_d\times D_d$ evaluation determinant is a nonzero polynomial in the stereographic
coordinates (exact rational instances of full rank: \texttt{certificates/corank\_exact.py c29},
$d=3,\dots,8$); the same genericity remark is in the proof
of P\'alv\"olgyi's Theorem~A.12 \cite{Palvolgyi}. Fewer generic points form a subset of a generic
spanning $D_d$-set, and subsets of Ramsey sets are Ramsey. These are corollaries of OpenAI's
Proposition~7.3; our part is noticing them. For $d=3$ this covers seven to nine generic points of $S^2$. They are not
two-group: four points of a M\"obius image of an algebraic set have an algebraic complex cross-ratio,
which four generic points do not, so each group has at most three points. On the circle ($D_2=5$) the
criterion gives at most five points, and every set of at most six cyclic points is two-group
(Remark~\ref{rem:six}). So non-two-group Ramsey sets off the circle are prior work, and
Theorem~\ref{thm:NTG} below is about the circle. With Proposition~\ref{prop:PD} applied on a great
circle, the consequence also shows that Ramsey-ness is not closed under specialization on $S^{d-1}$ for
$d\ge3$ (ledger SR-C7, SR-C16, SR-C29).
\end{remark}

On the circle, a set of at least four points is a M\"obius image of a set with algebraic coordinates if
and only if all its four-point cross-ratios are algebraic (cross-ratios are invariant under M\"obius maps
of the circle, i.e.\ $\mathrm{PGL}_2(\R)$ on the parameter; for the converse send three of its points to
$0,1,\infty$). A subset of such an image is again one, so a union of two such images with overlap is
also a union of two disjoint ones, and it suffices to test partitions.
Call a finite cyclic set \emph{two-group} if it is a union of two such images; two-group sets are Ramsey
(Theorem~\ref{thm:C2} with $d=2$). Throughout this section points are given by their stereographic
parameters (\S\ref{sec:prelim}, after a rational rotation off the pole), and a \emph{one-parameter
rational} set has parameters $r_i(t)$ with $r_i\in\Q(t)$ and $t$ transcendental, so that $F=\Q(t)$ and
$B=\Q(t)\otimes_\Q\Q(t)\subset\Q(T,U)$, with $x_i=r_i(T)$, $y_i=r_i(U)$. A non-constant element of $\Q(t)$
is transcendental, so cross-ratios of such a set are algebraic exactly when they are constant in $t$.

\begin{theorem}[Non-two-group Ramsey sets]\label{thm:NTG}
Let $t\in\R$ be transcendental. The cyclic sets with stereographic parameters
\[
 W_7=\{0,\tfrac12,\tfrac13,\ t,\ t+1,\ 2t,\ 1/t\},\qquad
 W_8=\{0,\tfrac12,\tfrac13,\tfrac23,\ t,\ t+1,\ 2t,\ 1/t\}
\]
are Ramsey and are not two-group. So is the radian $7$-gon $\{(\cos k\theta,\sin k\theta):0\le k\le6\}$
whenever $e^{i\theta}$ is transcendental (its parameters are $\tan(k\theta/2)\in\Q(\tan\tfrac\theta2)$).

\smallskip\noindent\emph{Status:} \st{CERT} plus glue. For each set an explicit $(2,2)$-form
$g(s,s')=\sum_{a,b\le2}g_{ab}s^as'^b$ with $g_{ab}\in\Q(T)[U]\subset B$ is printed in full by
\texttt{certificates/cert\_nontwogroup.py} (about 36\,s for $W_7,W_8$; with \texttt{--slow}, 37--47
minutes for the radian $7$-gon). The script checks exactly that $g(x_i,y_i)=0$ identically for every
point and that $g|_{U=T}=(s-s')^2$, and re-checks each certificate with a separate verification routine
in the same code; a non-builder re-check that parses the printed certificates and imports no code
from this project (sympy for $W_7,W_8$ in OpenAI's Cartesian form; python-flint for all three) also passes. The
glue is the circle dictionary of \S\ref{sec:prelim} and Lemma~0, exactly as for Theorem~\ref{thm:Ba};
it was reviewed (AI) as correct for these instances.
Not two-group (exact, same script): for $W_7$ and the radian $7$-gon none of the $35$ four-point
cross-ratios is constant in $t$, and every split of seven points into two parts has a part with at least
four points; for $W_8$ only the cross-ratio of the four rational points is constant, so a part with at
least four points lies in the four rational points, and two parts cover at most $4+3=7<8$ points; the
script also checks all splits. Scope \texttt{TRANSCENDENTAL\_PARAM}$(t)$, respectively
$(\tan\tfrac\theta2)$.

\smallskip\noindent\emph{What is not claimed.} The only structural claim is ``not two-group''. Ramsey
sets that are not subtransitive are already known (the Leader--Russell--Walters kites with
transcendental parameter, \cite[Cor.~7.5]{OAI172}), so we claim nothing about subtransitivity; we
do not claim that these sets are not subsoluble, which is open (\cite[Question~6.1]{Behague}); and we
do not claim that they are the first non-two-group Ramsey sets: on $S^2$ such sets follow from
OpenAI's Proposition~7.3 (Remark~\ref{rem:quadric}). The claim is about the circle.
\end{theorem}

\noindent So the two-group class of Theorem~\ref{thm:C2} is not the whole Ramsey class on the circle,
already at seven points. $W_8$ also shows that the phase diagram of Proposition~\ref{prop:PD} is a
statement about \emph{independent} points: four rational points plus four independent points are not
Ramsey, while four rational points plus four points on the rational curve $t\mapsto(t,t+1,2t,1/t)$ are.
The certificates were found by an exact decision procedure for one-parameter cyclic sets (kernel of the
evaluation map over $\Q(T,U)$, saturation at the diagonal, and a membership test for $(s-s')^2$;
\texttt{certificates/exact\_circle\_cert.py}), which reproduces five known exact verdicts as controls.

\medskip
\noindent\emph{The one-parameter criterion.} Let $A$ be a finite cyclic set whose coordinate field $F$ has
transcendence degree $1$ and is regular over $\Q$ ($\Q$ is algebraically closed in $F$), so that $B=F\otimes_\Q F$
is a domain; let $K$ be its fraction field and $\mathfrak p=\ker m$. Since $F$ has transcendence degree $1$,
$\mathcal O:=B_{\mathfrak p}$ is a discrete valuation ring with residue field $F$ (via $m$). Let
$V\subset K^9$ be the space of $(2,2)$-forms over $K$ that vanish at all pairs $(x_i,y_i)$, let
$L=V\cap\mathcal O^9$ (a saturated lattice, free of rank $\dim_KV$), and let $W=L/\mathfrak pL\subset F^9$ be its
reduction, an $F$-space of dimension $\dim_KV$.

\begin{lemma}[Limit criterion]\label{lem:limit}
\st{HAND}, reviewed (AI) as correct. $A$ is Ramsey if and only if $(s-s')^2\in W$.
\end{lemma}

\begin{proof}
If $A$ is Ramsey, the circle dictionary gives a $(2,2)$-form $f$ over $B$ vanishing at all $(x_i,y_i)$ with
$m(f)=\kappa(s-s')^2$, $\kappa\ne0$; then $f\in L$ and its reduction is $\kappa(s-s')^2$. Conversely, if
$g\in L$ reduces to $(s-s')^2$, write $g=f/c$ with $f\in B^9$ and $c\in B\setminus\mathfrak p$; then $f$
vanishes at all pairs and $m(f)=m(c)(s-s')^2$ with $m(c)\ne0$, so $A$ is Ramsey by the circle dictionary
and Lemma~0.
\end{proof}

\begin{theorem}[Rank criterion]\label{thm:rank}
Let $A$ be as above, with $5\le n=|A|\le8$. Fix $t\in F$ transcendental, let $\delta=d/dt$ on $F$, let
$s_i\in F$ be the stereographic parameters and $\delta_i=\delta(s_i)$, and let $N$ be the $n\times8$
matrix over $F$ with rows $(\delta_i,\ \delta_is_i,\ \delta_is_i^2,\ 1,\ s_i,\ s_i^2,\ s_i^3,\ s_i^4)$. If
$\operatorname{rank}_FN=n$, then $A$ is Ramsey.

\smallskip\noindent\emph{Status:} \st{HAND}, reviewed (AI; one non-author review, which found every step
correct and supplied the simplification used below). Scope: transcendence degree $1$, $F$ regular over
$\Q$.
\end{theorem}

\begin{proof}
$\dim_KV\ge9-n$, since $V$ is cut out by $n$ linear conditions on nine coefficients. Let $g_0\in W$. Reducing
$g(x_i,y_i)=0$ modulo $\mathfrak p$ gives $g_0(s_i,s_i)=0$; $g_0(s,s)$ is a polynomial of degree $\le4$ with
$n\ge5$ roots, so it vanishes and $g_0=(s-s')h_0$ with $h_0$ of bidegree $\le(1,1)$. The element
$\pi=1\otimes t-t\otimes1$ generates $\mathfrak p\mathcal O$, and $\mathcal O/\pi^2\mathcal O\cong F[\pi]/(\pi^2)$
with $F$ embedded as $F\otimes1$; in it $1\otimes s_i=s_i+\pi\delta_i$. Lift $g_0$ to $g\in L$ and write
$g\equiv g_0+\pi g_1$ modulo $\pi^2$. The coefficient of $\pi$ in $g(x_i,y_i)=0$ is
$\partial_{s'}g_0(s_i,s_i)\,\delta_i+g_1(s_i,s_i)=-h_0(s_i,s_i)\,\delta_i+g_1(s_i,s_i)=0$. So the coefficient
vector of $(-h_0(s,s),\,g_1(s,s))$, a quadratic and a quartic, lies in $\ker N$. Let
$\Phi:W\to F^3$ send $g_0$ to the coefficients of $h_0(s,s)$, which depend only on $g_0$. Then
$\Phi(W)\subset\pi_h(\ker N)$, the projection to the first three coordinates, and $\ker\Phi$ consists of the
$g_0=(s-s')h_0$ with $h_0$ vanishing on the diagonal, that is $\ker\Phi=W\cap F(s-s')^2$. If
$\operatorname{rank}N=n$, then $\dim\pi_h(\ker N)\le8-n<9-n\le\dim W$, so $\ker\Phi\ne0$, $(s-s')^2\in W$, and $A$ is
Ramsey by Lemma~\ref{lem:limit}.
\end{proof}

\noindent The criterion is consistent with the known failure: every point of P\'alv\"olgyi's heptagon
\cite{Palvolgyi} satisfies $\delta_i=-s_i(1+s_i^2)/(2r)$ in his parametrization, which puts two explicit
vectors in $\ker N$, so $\operatorname{rank}N\le6$ for every $r$ (\st{CERT}, an exact identity checked in
\texttt{certificates/rankN\_check.py}). On a seeded sample of $185$ one-parameter rational sets with
$6\le n\le8$ (including $40$ mirror-symmetric ones), $N$ has full rank, computed exactly at a rational
point (\texttt{certificates/rankN\_sweep.py}; the accompanying full-criterion passes are \st{PROBE}).

\begin{corollary}[The radian $8$-gon]\label{cor:radian8}
For every $\theta$ such that the eight points are distinct, $\{(\cos k\theta,\sin k\theta):0\le k\le7\}$ is
Ramsey.

\smallskip\noindent\emph{Status:} \st{HAND}, reviewed (AI), plus an exact rank computation: it
depends on Theorem~\ref{thm:rank} and on $\operatorname{rank}N=8$ at one rational value of $t$ (we
computed two). Scope \texttt{ANY}. This proves that a certificate exists; it does not exhibit one.
With Theorem~\ref{thm:NTG} (the radian $7$-gon) and Theorem~\ref{thm:D}, it shows that Theorem~\ref{thm:D}
is sharp: for transcendental $e^{i\theta}$ the radian progression is Ramsey with seven and eight points
and not with nine.
\end{corollary}

\begin{proof}
If $e^{i\theta}$ is algebraic, Theorem~\ref{thm:A}. Otherwise $t=\tan(\theta/2)$ is transcendental, $F=\Q(t)$
is regular over $\Q$, and $s_k=\tan(k\theta/2)$ and $\delta_k$ are explicit rational functions of $t$. The
matrix $N$ has rank $8$ at $t=2/7$ and at $t=-5/11$, computed exactly (\texttt{certificates/rankN\_check.py};
all denominators are nonzero there). Specialization does not raise rank, so $\operatorname{rank}N=8$ over
$\Q(t)$, and Theorem~\ref{thm:rank} applies.
\end{proof}

\noindent The reviewer re-verified the rank at three rational points and, by an
independent $\pi$-adic computation at $T=2/7$, found a kernel element whose reduction is proportional to
$(s-s')^2$.

\medskip
\noindent\emph{One rigid group.} For a cyclic set $A$ and $G\subset A$, say that $G$ has \emph{rational
cross-ratios} if all four-point cross-ratios within $G$ are rational.

\begin{theorem}[Rigid-group absorption]\label{thm:rigid}
Let $A$ be a finite subset of the unit circle whose coordinate field $F$ is regular over $\Q$, and suppose
some five points of $A$ have rational cross-ratios. Then $A$ is Ramsey if and only if $A=G\cup H$ with $G$
and $H$ each having rational cross-ratios, that is, if and only if $A$ is a union of two M\"obius images of
rational sets.

\smallskip\noindent\emph{Status:} \st{HAND}, reviewed (AI; the same non-author review as
Theorem~\ref{thm:rank}, which found every step correct). Scope: regular coordinate field
(\texttt{GENERIC\_OF} and \texttt{TRANSCENDENTAL\_PARAM} fields without algebraic irrationals).
The extension to algebraic seeds is open.
\end{theorem}

\begin{proof}
``If'' is Theorem~\ref{thm:C2} with $d=2$ (or Theorem~\ref{thm:Ba} with Theorem~\ref{thm:C}). For ``only
if'', let $f$ be a $(2,2)$-form over $B$ vanishing at all $(x_i,y_i)$ with $m(f)=\kappa(s-s')^2$,
$\kappa\ne0$ (circle dictionary); $B$ is a domain with fraction field $K$. We use that $c\otimes1=1\otimes c$
in $B$ forces $c\in\Q$. Fix three of the five given points, $g_1,g_2,g_3$, and let $G$ be the set of
$a\in A$ with $\mathrm{CR}(t_a,t_{g_1},t_{g_2},t_{g_3})\in\Q$; it contains the five points and has rational
cross-ratios. Let $\mu_G\in\PGL_2(K)$ send $x_{g_k}$ to $y_{g_k}$ ($k=1,2,3$). For $g\in G$ the cross-ratio
$c=\mathrm{CR}(t_g,t_{g_1},t_{g_2},t_{g_3})$ is rational, so $\mathrm{CR}(x_g,\dots)=c\otimes1=1\otimes c=\mathrm{CR}(y_g,\dots)$ and
$\mu_G(x_g)=y_g$. Let $g_G$ be the $(1,1)$-form of the graph of $\mu_G$ (irreducible).
(1) $f(s,\mu_G(s))$, cleared of denominators, has degree $\le4$ in $s$ and at least five roots $x_g$, so it
vanishes, and $f=g_G\,h$ in $K[s,s']$ with $h$ of bidegree $\le(1,1)$.
(2) For $a\in H:=A\setminus G$, $g_G(x_a,y_a)\neq0$: otherwise $\mu_G(x_a)=y_a$, so
$c=\mathrm{CR}(t_a,t_{g_1},t_{g_2},t_{g_3})$ satisfies $c\otimes1=1\otimes c$, hence $c\in\Q$ and $a\in G$. Since
$B$ is a domain, $h(x_a,y_a)=0$ for all $a\in H$.
(3) If $|H|\ge3$, $h$ is irreducible, because a reducible form of bidegree $\le(1,1)$ is a product of linear
forms in $s$ and in $s'$ and vanishes at no more than two pairs with distinct first and distinct second
coordinates. So $h$ is the graph of some $\mu\in\PGL_2(K)$ with $\mu(x_a)=y_a$ on $H$, and every four-point
cross-ratio $c$ of $H$ satisfies $c\otimes1=1\otimes c$, hence is rational. If $|H|\le3$ there is nothing to
prove.
\end{proof}

\noindent For five rational points plus $\mathrm{circ}(r_i(t))$, $r_i\in\Q(t)$, the theorem says: Ramsey iff at
most three of the $r_i$ are non-constant or all cross-ratios of the non-constant ones are constant in $t$.
On a seeded sample of $139$ such sets with four one-parameter points, the $119$ that are not two-group have
an exactly nonzero $9\times9$ evaluation determinant, hence are not Ramsey by the argument of
Theorem~\ref{thm:D} (\st{CERT} per instance), and the other $20$ are two-group, hence Ramsey
(\texttt{probes/c14\_rigid.py}).

\section{The non-Ramsey side}\label{sec:nonramsey}

\begin{lemma}[P\'alv\"olgyi's Lemma~A.1, re-derived]\label{lem:A1}
\st{HAND}. Let $A=\{a_i\}\subset\R^d$ be finite, $p_i=(1,a_i)$, and suppose there are real $\lambda_i$ and a
$\Q$-derivation $\partial$ of $\R$ (applied coordinatewise) such that $\sum_i\lambda_ip_ip_i^T=0$,
$\sum_i\lambda_i\partial a_i=0$, the matrix $\sum_i\lambda_i(\partial a_i)a_i^T$ is symmetric, and
$E:=\sum_i\lambda_i|\partial a_i|^2\ne0$. Then $A$ is not Ramsey.
\end{lemma}

\begin{proof}
Let $a_i'=v+Qa_i$ be a congruent copy in $\R^D$, $Q^TQ=I$. Then
$\partial a_i'=\partial v+(\partial Q)a_i+Q\,\partial a_i$. The moment conditions $\sum\lambda_i=0$,
$\sum\lambda_ia_i=0$, $\sum\lambda_ia_ia_i^T=0$ kill the weighted sum of $|\partial v+(\partial Q)a_i|^2$;
$\sum\lambda_i\partial a_i=0$ kills the cross term with $\partial v$; and
$2\sum_i\lambda_i\langle(\partial Q)a_i,Q\,\partial a_i\rangle=2\operatorname{tr}\bigl(K\sum_i
\lambda_i(\partial a_i)a_i^T\bigr)=0$, since $K=(\partial Q)^TQ$ is skew ($\partial(Q^TQ)=0$) and the other
factor is symmetric. So $z_i:=|\partial a_i'|^2$ satisfies $\sum_{i\ge2}\lambda_i(z_i-z_1)=E$. By
\cite[Lemma~15]{EGMRSS} there is a finite colouring $c$ of $\R$ with no monochromatic solution of this
equation; the colouring $x\mapsto c(|\partial x|^2)$ of $\R^D$ uses the same number of colours for every $D$
and has no monochromatic copy of $A$.
\end{proof}

\noindent The re-derivation was reviewed (AI) as correct, and the statement and proof of
\cite[Lemma~15]{EGMRSS} (pp.~349--350, proved in \S4, pp.~351--354, for any field) were read by a
reviewer; Rado's integer-variable Satz~XIII$'$ does not cover real coefficients. A derivation of the
coordinate field $F$ extends to $\R$ (transcendence-basis extension).

Consider $m$ rational points plus $j$ algebraically independent points on a circle, that is, $j$
points whose stereographic parameters are algebraically independent over $\Q$ (scope
\texttt{GENERIC\_OF}).

\begin{proposition}[PD, the circle phase diagram]\label{prop:PD}
Such a configuration is non-Ramsey if and only if $j\ge4$ and $m+j\ge7$.

\smallskip\noindent\emph{Status:} \st{HAND}, reviewed (AI) as correct. Ramsey side:
Theorem~\ref{thm:A} ($j=0$), Theorem~\ref{thm:Bb} ($j\le3$), Theorem~\ref{thm:Ba} ($m+j\le6$).
Non-Ramsey side: P\'alv\"olgyi's Theorems~A.3 and A.7 \cite{Palvolgyi}, plus closure of non-Ramsey sets
under supersets; the non-Ramsey half is his. Our write-up uses A.3, an instance of
Lemma~\ref{lem:A1} on the circle, in its stated general-triple form (and a rational rotation if a
rational point sits at the pole). His Theorem~A.3 is stated for an arbitrary real base triple
$t_1,t_2,t_3$ with $t_4,\dots,t_7$ algebraically independent over $\Q(t_1,t_2,t_3)$, so it covers every
cell $m+j=7$, $j\ge4$, including $m<3$ (take the base to be the $m$ rational points and $3-m$
independent ones). His Theorem~A.7 (seven points, every six of which have three algebraically
independent projective moduli, are not Ramsey) also covers $(0,7),(1,6),(2,5)$ and $(3,4)$: a
six-point subset has at most three rational points, and sending three of its points, all its rational
ones among them, to $\infty,0,1$ leaves three algebraically independent moduli. A.7 rests on his A.5
and A.6 as well, so we cite A.3. (This check of A.3's and A.7's hypotheses was reviewed (AI), with an
exact Jacobian-rank check; A.3 and A.7 are in his appendix, which he labels as unchecked, and our route
rests on our re-derivation of Lemma~A.1.)
\end{proposition}

\noindent\emph{Exact instance (J2-F7).} \st{CERT}, scope
\texttt{GENERIC\_OF}(3 rational $+$ 4 independent). For the rational parameters $0,\tfrac12,\tfrac13$
and four independent parameters $T_1,\dots,T_4$, the scalar certificate with the derivation $\partial=\sum_mg(T_m)\,\partial/\partial T_m$,
$g(t)=t(t-\tfrac12)(t-\tfrac13)$, satisfies the hypotheses of Lemma~\ref{lem:A1} exactly, with $E=4$
(\texttt{certificates/dual\_cert.py}), so the configuration is non-Ramsey. A concrete member is
$(T_m)=(e^{\sqrt2},e^{\sqrt3},e^{\sqrt5},e^{\sqrt7})$, by Lindemann--Weierstrass. This is an instance of
P\'alv\"olgyi's Theorem~A.3, which is prior.

\noindent\emph{Credit summary for the circle.} Seven points: P\'alv\"olgyi's heptagon \cite{Palvolgyi}
is the first non-Ramsey cyclic set; Theorem~\ref{thm:Ba} (prior, \cite{JDJones}) gives the minimum. The
non-Ramsey half of Proposition~\ref{prop:PD} is P\'alv\"olgyi's (A.3, A.7); the Ramsey half combines
prior work with Theorem~\ref{thm:Bb}.

\begin{theorem}[NR-S23]\label{thm:NR}
The following configurations are not Ramsey: on $S^2$, the first $m$ of the rational points $R_2$
below plus $j$ points with algebraically independent stereographic coordinates, for
$(m,j)=(6,4),(5,5),(4,6),(3,7)$, and ten points with all coordinates independent; on $S^3$, the
first $m$ of the rational points $R_3$ below plus $j$ independent points, for $(m,j)=(11,4),(10,5)$.
Hence so is every superset, which covers all sixteen failing cells of the $S^2$ scan and both
failing threshold cells on $S^3$.

The rational points, as stereographic coordinates $u$ from the pole $(0,\dots,0,1)$ (the point is
$(2u,|u|^2-1)/(1+|u|^2)$), in order:
\begin{align*}
R_2:\ &(\tfrac59,\tfrac58),\ (\tfrac74,-\tfrac49),\ (2,-\tfrac34),\ (0,-\tfrac79),\ (-\tfrac87,\tfrac53),\
(-1,-7);\\
R_3:\ &(-\tfrac15,-1,-\tfrac54),\ (-\tfrac43,-\tfrac79,-\tfrac35),\ (-\tfrac97,-5,-\tfrac13),\
(-\tfrac75,\tfrac57,-1),\ (\tfrac12,\tfrac23,\tfrac97),\ (\tfrac27,\tfrac12,\tfrac25),\\
&(\tfrac59,-\tfrac58,\tfrac59),\ (-\tfrac45,-1,\tfrac76),\ (-\tfrac16,1,-\tfrac17),\
(-\tfrac58,2,\tfrac94),\ (-\tfrac12,-\tfrac52,\tfrac12).
\end{align*}
For example the first point of $R_2$ is $(5760,6480,-1559)/8809\in S^2$.

\smallskip\noindent\emph{Credit.} The cell of ten points with all coordinates independent is the case
$d=3$ of P\'alv\"olgyi's Theorem~A.12 \cite{Palvolgyi} (appendix, labelled unchecked by its author),
which is prior; what is ours there is the exact certificate.

\smallskip\noindent\emph{Status:} \st{CERT} identities plus glue, reviewed (AI) as correct. The
scalar certificates are exact (\texttt{certificates/scalar\_exact.py}; replay byte-identical): cofactor
vectors give $\lambda$ and $\partial$ identically in the free coordinates, and at one rational specialization
$\tau$ both systems have full rank and $E(\tau)\ne0$, so $E\ne0$ in $F$. Two valid routes, in every
dimension: Lemma~\ref{lem:A1} (no Lean, no J2), or the 2-jet dual and Proposition~\ref{prop:J2}. A code
check: the script verifies $\tilde\lambda_i(\tau)\ne0$ for every point that carries the derivation; for the
other points whose weight is not identically zero the released script checks the same, and passes in
all seven cells. Scope:
\texttt{GENERIC\_OF}(independent free coordinates, with the rational points $R_2$, $R_3$ above).
\end{theorem}

\noindent\emph{$S^4$.} The analogous threshold cell $(17,4)$ is not yet certified: the $(P)$ system has a
second, identically vanishing row relation (an infinitesimal M\"obius rotation fixing the 2-sphere through
the four free points), so one more row must be dropped. With that change the reviewer's scratch run gives
an exact $E(\tau)\neq0$; that computation is not part of this release. A scalar certificate exists modulo
$p$ \st{PROBE}. Since version~1.1, ledger entry SR-C30 certifies the cell $(N(d)-4,4)$ exactly (the
rank computation absorbs the dependent rows) for seeded rational points in general position for $d=3,\dots,7$, which includes a
cell $(17,4)$ on $S^4$ (\texttt{certificates/corank\_exact.py c30}; glue reviewed once by an AI agent).

\section{A second-order necessary test}\label{sec:jet}

Let $F$ have transcendence degree $k$ over $\Q$, with commuting derivations $D_1,\dots,D_k$ dual to a
transcendence basis. Expanding condition (1) to second order along the diagonal of $B$ gives a linear
system in $c,P_j,Q_{jl}$ (the ``2-jet system''): $p_i^TP_jp_i=0$ and
$-c\langle D_ja_i,D_la_i\rangle+p_i^TP_jD_lp_i+p_i^TP_lD_jp_i+p_i^TQ_{jl}p_i=0$.

\begin{proposition}[J2]\label{prop:J2}
\st{HAND}, scope \texttt{ANY}. If $A$ is Ramsey, the 2-jet system is solvable with $c=1$. Proof: the
Taylor map $f\otimes g\mapsto f\cdot\sum_\alpha D^\alpha g\,e^\alpha/\alpha!$ is a ring map
$B\to F[e]/(e)^3$ compatible with $m$; apply it to a symmetrized solution. Equivalently, the 2-jet test
is the field criterion over $B/\mathfrak m^3$, where $\mathfrak m=\ker m$, so it does not depend on the
transcendence basis. This proof was found and written by an AI cold reviewer of our original derivation,
who confirmed the order-1 and order-2 coefficients, signs and normalization; there has been no human
review.
\end{proposition}

\noindent\emph{Dual form (corrected).} By the Fredholm alternative over $F$, non-solvability is
witnessed by symmetric $k\times k$ weights $\Lambda_i$ with $\sum_i(\Lambda_i)_{jl}p_ip_i^T=0$, with
$\sum_i\sum_l(\Lambda_i)_{jl}\,p_i(D_lp_i)^T\in\mathrm{span}_F\{p_ip_i^T\}$ for each $j$ as an equality of
full matrices, and with $\sum_i\operatorname{tr}(\Lambda_iG_i)\neq0$, where $G_i$ is the Gram matrix of the
$D_ja_i$. The version with ``$=0$'' in the middle condition, as first written, is sound but incomplete: it
misses real failures (exact check on the J2-F7 configuration). P\'alv\"olgyi's conditions
(Lemma~\ref{lem:A1}) are the case $\Lambda_i=\lambda_iC$, and his appendix reduces a general $C$ to a
single derivation. An earlier claim of ours that matrix weights are strictly stronger is
\textbf{withdrawn}: its witness has a scalar certificate (J2-F7).

\noindent\emph{Evidence for sufficiency.} \st{PROBE} For one parameter ($k=1$) the full criterion is
computable, because the local ring at the diagonal is a discrete valuation ring. On 702 real
one-parameter configurations where both tests gave a verdict (635 pass/pass, 67 fail/fail), the 2-jet
test and the full criterion agree in every case, at two primes near $2^{31}$. ``Real'' means an open
parameter interval on which every radicand is positive, decided exactly; an earlier sample that included
non-real configurations is not used. These are random specializations modulo a prime; they are evidence
for Conjecture~\ref{conj:C1}, not theorems. The test also reproduces every known verdict
\cite{OAI172,Palvolgyi}.

\noindent\emph{Scalar certificates.} \st{PROBE} In every tested failing cell (all sixteen on $S^2$, both
$S^3$ threshold cells, the $S^4$ threshold, and ten generic points on $S^2$) a scalar single-derivation
certificate exists, at two primes; seven of them are exact (Theorem~\ref{thm:NR}). A random fixed
direction is not a test for this: it misses certificates that exist. \emph{Re-scoped in version~1.1:}
this agreement is automatic and is not evidence about ledger entries SR-C11 and SR-C12 (whether a 2-jet
failure without a scalar certificate exists, and the working hypothesis that none exists on
configurations of rational plus independent points). Every cell above has quadric corank one (one
relation $\sum_i\lambda_ip_ip_i^T=0$) or contains a failing cell of corank one, and at corank one the
dual space of the 2-jet system consists of weights $\Lambda_i=\lambda_iC$, so a 2-jet failure there
gives a scalar certificate once $C$ is reduced to a single derivation (\st{HAND}, reviewed twice by
non-author AI agents; ledger SR-C31(a); written out in Remark~\ref{rem:corankjet}). 2-jet failures with no scalar certificate exist at quadric corank
three and four on $S^2$ (ledger SR-C32, proved for four seeded families), which answers SR-C11.

\section{First-order rigidity}\label{sec:rigidity}

\noindent\emph{Status of this section.} \st{HAND}: written proof in this section (added in v1.1.1),
independently reviewed twice by non-author AI agents (the first review reconstructed the argument; a
second, independent review found every step sound; \S\ref{sec:disclosure}). Scope \texttt{ANY}. Version~1.1.0 stated first-order rigidity (ledger SR-C34) and the
finite-type reformulation (ledger SR-C33(a)--(c)) with a summary of the proof only; the complete argument
is written out here. It follows the reconstruction by the non-author AI reviewer of version~1.1.0
(\S\ref{sec:disclosure}). It rests on OpenAI's Theorem~1.1 and its Lean proof (\S\ref{sec:prelim}); the
rest is standard commutative algebra over $F\otimes_\Q F$, and each standard fact is named where it is
used.

\medskip\noindent\emph{Setting.} Throughout this section $A=\{a_1,\dots,a_s\}\subset\R^d$ consists of
$s\ge2$ distinct points with $|a_i|=1$ (the unit sphere centred at $0$), $d\ge1$. Ramsey-ness is invariant
under similarities, so a set on another sphere is first moved to this one; the statements below concern
the coordinates of the moved set, whose coordinate field is in general a different field. Let
$p_i=(1,a_i)^T$, $F=\Q(a_{i\alpha})$, $B=F\otimes_\Q F$, $m:B\to F$ the multiplication map,
$\mathfrak m=\ker m$, and $\tau$ the flip $x\otimes y\mapsto y\otimes x$. Put $M=(d+1)^2$, identify
$(d+1)\times(d+1)$ matrices over any ring with $M$-tuples, and let
\[
 \Phi(P)_i=(p_i\otimes1)^TP(1\otimes p_i)=\sum_{\alpha,\beta}(p_{i\alpha}\otimes p_{i\beta})P_{\alpha\beta}
 \qquad(1\le i\le s),
\]
a $B$-linear map $B^M\to B^s$, given by the same formula over every $B$-algebra. Let
$S=\mathrm{diag}(-1,I_d)$ and $b:=\Phi(S\otimes1)\in B^s$. Applying $m$ to the entries of $\Phi$ gives
$\bar\Phi:\Mat_{d+1}(F)\to F^s$, $X\mapsto(p_i^TXp_i)_i$, whose rows are the matrices $p_ip_i^T$. Let
$\rho$ be its rank (the quadric rank), $\kappa=s-\rho$ the quadric corank, and
$\mathrm{Rel}=\{\lambda\in F^s:\sum_i\lambda_ip_ip_i^T=0\}$, of dimension $\kappa$. Let $k$ be the
transcendence degree of $F$ over $\Q$, $t_1,\dots,t_k\in F$ a transcendence basis, and $D_1,\dots,D_k$
the derivations of $F$ with $D_jt_l=\delta_{jl}$. They exist, are unique and commute, because $F$ is a finite
separable extension of $\Q(t)$: a derivation of $\Q(t)$ extends uniquely along a separable algebraic
extension (if $\mu$ is the minimal polynomial of $g$, then $0=D(\mu(g))=\mu^D(g)+\mu'(g)Dg$, where $\mu^D$ applies $D$ to the coefficients, with
$\mu'(g)\ne0$), and $[D_j,D_l]$ is a derivation vanishing on $\Q(t)$. They form an $F$-basis of
$\mathrm{Der}_\Q(F)$, since $\partial-\sum_j\partial(t_j)D_j$ vanishes on $\Q(t)$, hence on $F$. For
$g\in F$ write $\delta g=1\otimes g-g\otimes1\in\mathfrak m$. Nothing is assumed about $F$ beyond finite
generation over $\Q$: $B$ need not be a domain, and $F$ need not be regular over $\Q$.

\begin{lemma}[The local ring at the diagonal]\label{lem:local}
\begin{enumerate}
\item[(a)] $B$ is Noetherian.
\item[(b)] $\mathfrak m$ is generated, as a module over $F\otimes1$, by the elements $\delta g$, and
$\delta g\equiv\sum_jD_j(g)\,\delta t_j$ modulo $\mathfrak m^2$.
\item[(c)] Let $\sigma(g)=\sum_{|\alpha|\le2}D^\alpha g\,e^\alpha/\alpha!\in F[e_1,\dots,e_k]/(e)^3$. The Taylor
map $T(f\otimes g)=f\,\sigma(g)$ induces an $F\otimes1$-linear ring isomorphism
$B/\mathfrak m^3\to F[e]/(e)^3$, with $T(\mathfrak m^r/\mathfrak m^3)=(e)^r/(e)^3$ and
$T(\delta g)\equiv\sum_je_jD_jg$ modulo $(e)^2$. In particular $\mathfrak m/\mathfrak m^2$ has $F$-basis
$\delta t_1,\dots,\delta t_k$, and $\mathfrak m^2/\mathfrak m^3$ has $F$-basis $\delta t_j\delta t_l$
($j\le l$).
\item[(d)] $\mathcal O:=B_{\mathfrak m}$ is a Noetherian local ring with maximal ideal $\mathfrak m\mathcal O$
and residue field $F$. For $n\ge1$ the natural map $B/\mathfrak m^n\to\mathcal O/\mathfrak m^n\mathcal O$ is
an isomorphism carrying $\mathfrak m^r/\mathfrak m^n$ onto $\mathfrak m^r\mathcal O/\mathfrak m^n\mathcal O$
for every $0\le r\le n$; in particular $\mathfrak m\mathcal O\cap B=\mathfrak m$.
\item[(e)] If $k=1$, $\mathcal O$ is a discrete valuation ring with uniformizer $\pi=\delta t_1$.
\end{enumerate}
\end{lemma}

\begin{proof}
(a) $\Q(t)\otimes_\Q F$ is the localization of $\Q[t]\otimes_\Q F\cong F[X_1,\dots,X_k]$ (a polynomial ring, the $t_j$ being algebraically independent) at the nonzero elements
of $\Q[t]\otimes1$, so it is Noetherian by Hilbert's basis theorem. $B=F\otimes_{\Q(t)}(\Q(t)\otimes_\Q F)$
is a free module of rank $[F:\Q(t)]<\infty$ over it, and every ideal of $B$ is a submodule of this finitely
generated module, hence finitely generated.

(b) If $x=\sum_rf_r\otimes g_r\in\mathfrak m$, then $\sum_rf_rg_r=0$ and
$x=\sum_r(f_r\otimes1)\,\delta g_r$. The quotient $\mathfrak m/\mathfrak m^2$ is a vector space over
$B/\mathfrak m=F$, on which $f\otimes1$ and $1\otimes f$ act alike, and
$\delta(gh)=(1\otimes g)\,\delta h+(h\otimes1)\,\delta g$, so $g\mapsto\delta g$ is a $\Q$-derivation
$F\to\mathfrak m/\mathfrak m^2$. Then $g\mapsto\delta g-\sum_jD_j(g)\,\delta t_j$ is a derivation vanishing
on $t_1,\dots,t_k$, hence on $\Q(t)$ and on $F$, as in the setting.

(c) The $D_j$ commute, so the Leibniz rule
$D^\alpha(gh)=\sum_{\beta\le\alpha}\binom\alpha\beta D^\beta g\,D^{\alpha-\beta}h$ shows that $\sigma$ is a ring
homomorphism modulo $(e)^3$; $T$ is the ring homomorphism $B\to F[e]/(e)^3$ induced by the two
$\Q$-algebra maps $f\mapsto f$ and $\sigma$ (this is the Taylor map of Proposition~\ref{prop:J2}). Setting $e=0$
recovers $m$, so $T(\mathfrak m)\subset(e)$ and $T(\mathfrak m^3)=0$. Since $T(f\otimes1)=f$ and
$T(\delta t_j)=e_j$, $T$ is onto. By (b), $\mathfrak m/\mathfrak m^2$ is spanned over $F$ by the $k$
classes $\delta t_j$, and $\mathfrak m^2/\mathfrak m^3$ by the products $\delta t_j\delta t_l$; so
$\dim_FB/\mathfrak m^3\le1+k+k(k+1)/2=\dim_FF[e]/(e)^3$, and the onto $F$-linear map is bijective. The
formula for $T(\delta g)=\sigma(g)-g$ is immediate.

(d) A localization of a Noetherian ring is Noetherian. Every prime ideal containing $\mathfrak m^n$
contains $\mathfrak m$, so $B/\mathfrak m^n$ is local and every $c\notin\mathfrak m$ is a unit modulo
$\mathfrak m^n$. Injectivity: if $x\in B$ equals $y/c$ in $\mathcal O$ with $y\in\mathfrak m^n$,
$c\notin\mathfrak m$, then $u(cx-y)=0$ for some $u\notin\mathfrak m$; as $uc$ is a unit modulo
$\mathfrak m^n$, $x\in\mathfrak m^n$. Surjectivity: $y/c\equiv yc'$ modulo $\mathfrak m^n\mathcal O$, where
$cc'\equiv1$ modulo $\mathfrak m^n$. Every element of $\mathfrak m^r\mathcal O$ is $y/c$ with $y\in\mathfrak m^r$,
and then $yc'\in\mathfrak m^r$, so the image of $\mathfrak m^r/\mathfrak m^n$ is all of
$\mathfrak m^r\mathcal O/\mathfrak m^n\mathcal O$. The case $n=1$ gives the last claim.

(e) $B=F\otimes_{\Q(t_1)}\bigl(\Q(t_1)\otimes_\Q\Q(t_1)\bigr)\otimes_{\Q(t_1)}F$ is a free module over
$R_0=\Q(t_1)\otimes_\Q\Q(t_1)$, a localization of the domain $\Q[T,U]$, in which $\pi=U-T\ne0$. So
multiplication by $\pi$ is injective on $B$, hence on $\mathcal O$ (localization is exact). By (c) and (d),
$\mathfrak m\mathcal O=\pi\mathcal O+\mathfrak m^2\mathcal O$, so $\mathfrak m\mathcal O=\pi\mathcal O$ by
Nakayama's lemma. By Krull's intersection theorem $\bigcap_n\pi^n\mathcal O=0$, so every nonzero
$x\in\mathcal O$ is $\pi^nu$ with $u\notin\pi\mathcal O$, a unit; since $\pi$ is a nonzerodivisor, products
of nonzero elements are nonzero. So $\mathcal O$ is a domain whose nonzero ideals are the $\pi^n\mathcal O$.
\end{proof}

\noindent(One can show that $\mathcal O$ is regular of dimension $k$; nothing below uses this.)

\begin{lemma}[The obstruction]\label{lem:obstruction}
$b_i=-\tfrac12\sum_{j=1}^d(\delta a_{ij})^2$ for every $i$. Hence $b\in\mathfrak m^2B^s$, and
$T(b_i)=-\tfrac12\,e^TG_ie$, where $G_i$ is the Gram matrix $(\langle D_ja_i,D_la_i\rangle)_{j,l}$.
\end{lemma}

\begin{proof}
$b_i=\sum_{\alpha,\beta}S_{\alpha\beta}\,p_{i\alpha}\otimes p_{i\beta}=-1\otimes1+\sum_ja_{ij}\otimes a_{ij}$,
while $(\delta a_{ij})^2=1\otimes a_{ij}^2-2a_{ij}\otimes a_{ij}+a_{ij}^2\otimes1$ and
$\sum_ja_{ij}^2=1$ in $F$. The last claim follows from Lemma~\ref{lem:local}(c):
$T(\delta a_{ij})^2\equiv(\sum_le_lD_la_{ij})^2$ modulo $(e)^3$.
\end{proof}

\noindent By the \emph{field criterion over $B/\mathfrak m^n$} ($n\ge1$) we mean a matrix $P$ over
$B/\mathfrak m^n$ with $\Phi(P)=0$ in $(B/\mathfrak m^n)^s$ and spatial block of $m(P)$ equal to $I_d$ ($m$
factors through $B/\mathfrak m^n$). For $n=3$ this is the 2-jet test (Proposition~\ref{prop:J2}).

\begin{proposition}[Finite type; ledger SR-C33(a)]\label{prop:finitetype}
\begin{enumerate}
\item[(i)] If $b\in\Phi(\mathfrak m\mathcal O^M)$, then $A$ is Ramsey. (No spanning hypothesis.)
\item[(ii)] If $A$ affinely spans $\R^d$ and is Ramsey, then $b\in\Phi(\mathfrak mB^M)$.
\item[(iii)] If $A$ affinely spans $\R^d$ and $n\ge1$, the field criterion over $B/\mathfrak m^n$ is
solvable if and only if $b\in\Phi(\mathfrak mB^M)+\mathfrak m^nB^s$.
\item[(iv)] There is $n_0=n_0(A)$ such that for every $n\ge n_0$:
$b\in\Phi(\mathfrak m\mathcal O^M)$ if and only if $b\in\Phi(\mathfrak mB^M)+\mathfrak m^nB^s$. So a set
that affinely spans $\R^d$ is Ramsey if and only if the field criterion over $B/\mathfrak m^n$ is solvable,
for any $n\ge n_0(A)$.
\end{enumerate}
The proof gives no bound on $n_0(A)$. Conjecture~\ref{conj:C1} says that $n_0=3$ works for every spanning
set.
\end{proposition}

\begin{proof}
(i) Write $b=\Phi(Q)$ with $Q\in(\mathfrak m\mathcal O)^M$, and $Q=q/c$ with $q\in B^M$ and
$c\in B\setminus\mathfrak m$ a common denominator. Each $q_{\alpha\beta}=cQ_{\alpha\beta}$ lies in
$\mathfrak m\mathcal O\cap B=\mathfrak m$ (Lemma~\ref{lem:local}(d)). The equality $cb=\Phi(q)$ in
$\mathcal O^s$ means $u(cb-\Phi(q))=0$ in $B^s$ for some $u\notin\mathfrak m$ (a product of one such element
per coordinate). Put $P=uc\,(S\otimes1)-uq$. Then $\Phi(P)=u(cb-\Phi(q))=0$, and $m(P)=m(uc)\,S$ since
$m(q)=0$; here $m(uc)\ne0$ because $\mathfrak m$ is prime. So the spatial block of $m(P)$ is $m(uc)I_d$ with
$m(uc)\in F^\times$, and by Lemma~0 $(m(uc)^{-1}\otimes1)P$ satisfies OpenAI's field criterion. Their Lean
sufficiency \texttt{fieldCriterion\_\allowbreak{}ramsey} needs only $s\ge2$, $d\ge1$ and distinct points
(\S\ref{sec:prelim}), so $A$ is Ramsey.

(ii) By Theorem~1.1 (Lean: \texttt{ramsey\_\allowbreak{}realCriterion}, then
\texttt{realCriterion\_\allowbreak{}fieldCriterion}) there is $P$ over $B$ with $\Phi(P)=0$ and spatial
block of $m(P)$ equal to $I_d$. Let $P^*_{\alpha\beta}=\tau(P_{\beta\alpha})$. Since $\tau$ is a ring
automorphism with $m\circ\tau=m$, $\Phi(P^*)_i=\tau(\Phi(P)_i)=0$ and $m(P^*)=m(P)^T$. So
$P'=\tfrac12(P+P^*)$ has $\Phi(P')=0$ and $m(P')$ symmetric with spatial block $I_d$. Applying $m$ to
$\Phi(P')_i=0$ gives $p_i^Tm(P')p_i=0$: the function $x\mapsto(1,x^T)\,m(P')\,(1,x^T)^T=|x|^2+2v\cdot x+c_0$
vanishes on $A$, and so does $|x|^2-1$. Their difference is affine and vanishes on a set that affinely
spans $\R^d$, so $v=0$, $c_0=-1$ and $m(P')=S$. Then $Q:=S\otimes1-P'$ has $m(Q)=0$, so its entries lie in
$\mathfrak m$, and $\Phi(Q)=b$.

(iii) If $b=\Phi(Q)+r$ with $Q\in\mathfrak m^M$ and $r\in\mathfrak m^nB^s$, then $S\otimes1-Q$ solves the
criterion over $B/\mathfrak m^n$. Conversely, $\tau(\mathfrak m)=\mathfrak m$, so $\tau$ acts on
$B/\mathfrak m^n$ and the symmetrization and the spanning argument of (ii) apply verbatim to a solution
over $B/\mathfrak m^n$; lift the resulting solution with $m(P')=S$ to $\hat P$ over $B$. Then
$Q=S\otimes1-\hat P$ has entries in $\mathfrak m$ and $b-\Phi(Q)=\Phi(\hat P)\in\mathfrak m^nB^s$.

(iv) Let $K=\Phi(\mathfrak m\mathcal O^M)$. By Lemma~\ref{lem:local}(d) (reduce modulo $\mathfrak m^n$ and
lift), $b\in K+\mathfrak m^n\mathcal O^s$ if and only if $b\in\Phi(\mathfrak mB^M)+\mathfrak m^nB^s$. If
$b\in K$, take $n_0=1$. Otherwise the image of $b$ in the finitely generated $\mathcal O$-module
$\mathcal O^s/K$ is nonzero, and Krull's intersection theorem ($\bigcap_n\mathfrak m^nN=0$ for every finitely
generated module $N$ over a Noetherian local ring) gives $n_0$ with $b\notin K+\mathfrak m^{n_0}\mathcal O^s$,
hence $b\notin K+\mathfrak m^n\mathcal O^s$ for all $n\ge n_0$. The last sentence combines this with (i),
(ii) and (iii).
\end{proof}

\begin{lemma}[Smith form at the diagonal]\label{lem:smith}
There are invertible matrices $U$ ($s\times s$) and $V$ ($M\times M$) over $\mathcal O$ with
\[
 U^{-1}\Phi V^{-1}=\begin{pmatrix}I_\rho&0\\0&\Psi\end{pmatrix},
\]
where $\Psi$ is a $\kappa\times(M-\rho)$ matrix with entries in $\mathfrak m\mathcal O$. Let $\beta\in\mathcal O^\kappa$
be the last $\kappa$ coordinates of $U^{-1}b$, and $\operatorname{Im}\Psi\subset\mathcal O^\kappa$ the
column span of $\Psi$. Then $\beta\in\mathfrak m^2\mathcal O^\kappa$, and $b\in\Phi(\mathfrak m\mathcal O^M)$ if
and only if $\beta\in\mathfrak m\operatorname{Im}\Psi$. The last $\kappa$ rows of $U_0^{-1}$, where $U_0$ is the
reduction of $U$ modulo $\mathfrak m$, form a basis of $\mathrm{Rel}$.
\end{lemma}

\begin{proof}
Invertible row and column operations over $\mathcal O$ do not change the rank of the reduction modulo
$\mathfrak m$. If some entry of a matrix is a unit, move it to position $(1,1)$, scale, and clear its row and
column; the remaining block has reduction of rank one less. Repeat until every entry of the remaining
block lies in $\mathfrak m\mathcal O$; this happens after exactly $\rho$ steps. Since $V$ preserves
$\mathfrak m\mathcal O^M$,
$\Phi(\mathfrak m\mathcal O^M)=U\bigl(\mathfrak m\mathcal O^\rho\oplus\mathfrak m\operatorname{Im}\Psi\bigr)$.
By Lemma~\ref{lem:obstruction}, $U^{-1}b\in\mathfrak m^2\mathcal O^s$, so its first $\rho$ coordinates lie in
$\mathfrak m\mathcal O$ automatically, which gives the two claims about $\beta$. Finally
$U_0^{-1}\bar\Phi=\begin{psmallmatrix}I_\rho&0\\0&0\end{psmallmatrix}V_0$, so the last $\kappa$ rows of the
invertible matrix $U_0^{-1}$ are independent vectors $\ell$ with $\ell\bar\Phi=0$, that is, with
$\sum_i\ell_ip_ip_i^T=0$; there are $\kappa=\dim\mathrm{Rel}$ of them.
\end{proof}

\begin{definition}[The 2-jet dual space]\label{def:Delta}
$\Delta(A)$ is the $F$-space of tuples $\Lambda=(\Lambda_1,\dots,\Lambda_s)$ of symmetric $k\times k$
matrices over $F$ such that
\begin{enumerate}
\item[(D1)] $\sum_i(\Lambda_i)_{jl}\,p_ip_i^T=0$ for all $j,l$, and
\item[(D2)] $X_j:=\sum_i\sum_l(\Lambda_i)_{jl}\,p_i(D_lp_i)^T\in\mathrm{span}_F\{p_ip_i^T\}$ for each $j$,
as an equality of full matrices.
\end{enumerate}
These are the first two conditions of the corrected dual form in \S\ref{sec:jet}; the trace condition is not
imposed. For $\lambda\in F^s$ let $\Pi_\lambda:F^k\to\mathrm{Alt}_{d+1}(F)$ (antisymmetric matrices, of
dimension $d(d+1)/2=\dim\mathfrak{so}(1,d)$) be
$\Pi_\lambda(\xi)=\operatorname{asym}\bigl(\sum_i\lambda_ip_i(\sum_l\xi_lD_lp_i)^T\bigr)$, where
$\operatorname{asym}(X)=\tfrac12(X-X^T)$ and $\operatorname{sym}(X)=\tfrac12(X+X^T)$.
\end{definition}

\begin{lemma}[Symmetric parts are automatic]\label{lem:Pi}
\begin{enumerate}
\item[(a)] For $\lambda\in\mathrm{Rel}$ and any $\Q$-derivation $\partial$ of $F$, the matrix
$Y=\sum_i\lambda_ip_i(\partial p_i)^T$ satisfies $Y+Y^T=-\sum_i(\partial\lambda_i)\,p_ip_i^T$; so
$\operatorname{sym}(Y)\in\mathrm{span}_F\{p_ip_i^T\}$.
\item[(b)] Let $\lambda^1,\dots,\lambda^\kappa$ be a basis of $\mathrm{Rel}$. Writing
$\Lambda_i=\sum_r\lambda^r_iC^r$ identifies $\Delta(A)$ with the space of tuples $(C^1,\dots,C^\kappa)$ of
symmetric $k\times k$ matrices such that $\sum_r\Pi_{\lambda^r}(C^r\xi)=0$ for all $\xi\in F^k$. In
particular $\Delta(A)=0$ if $\kappa=0$, and $\Delta(A)\cong\mathrm{Sym}^2(\ker\Pi_\lambda)$ if
$\mathrm{Rel}=F\lambda$. The same description holds with any $F$-basis of $\mathrm{Der}_\Q(F)$ in place of
$D_1,\dots,D_k$, up to the isomorphism $C\mapsto g^TCg$ for the change of basis $g$.
\end{enumerate}
\end{lemma}

\begin{proof}
(a) Apply $\partial$ entrywise to $\sum_i\lambda_ip_ip_i^T=0$:
$\sum_i(\partial\lambda_i)p_ip_i^T+\sum_i\lambda_i\bigl((\partial p_i)p_i^T+p_i(\partial p_i)^T\bigr)=0$.

(b) (D1) says that for each $(j,l)$ the vector $((\Lambda_i)_{jl})_i$ lies in $\mathrm{Rel}$, so
$\Lambda_i=\sum_r\lambda^r_iC^r$ for unique symmetric $C^r$. Then
$X_j=\sum_r\sum_lC^r_{jl}Y^r_l$ with $Y^r_l=\sum_i\lambda^r_ip_i(D_lp_i)^T$. By (a),
$\operatorname{sym}(X_j)\in\mathrm{span}_F\{p_ip_i^T\}$, and that span consists of symmetric matrices, so
(D2) holds if and only if $\operatorname{asym}(X_j)=\sum_r\Pi_{\lambda^r}(C^re_j)=0$ ($C^r$ is symmetric).
For $\kappa=1$ the condition says that every column of $C$ lies in $\ker\Pi_\lambda$. For another basis
$D'_l=\sum_jg_{lj}D_j$, $\Pi'_\lambda(\xi)=\Pi_\lambda(g^T\xi)$, and $C'\mapsto g^TC'g$ carries one
solution space onto the other.
\end{proof}

\noindent Let $L_2:(\mathfrak m/\mathfrak m^3)^M\to(\mathfrak m/\mathfrak m^3)^s$ be the $F$-linear map
induced by $\Phi$ (the 2-jet system), and $I_2=\operatorname{Im}L_2\cap(\mathfrak m^2/\mathfrak m^3)^s$. By
Lemma~\ref{lem:local}(c), $T$ identifies $(\mathfrak m^2/\mathfrak m^3)^s$ with $s$-tuples of quadratic forms
$e^TX_ie$, $X_i$ symmetric, and $\langle\Lambda,X\rangle=\sum_i\operatorname{tr}(\Lambda_iX_i)$ is a
perfect pairing with tuples of symmetric $k\times k$ matrices.

\begin{lemma}[Duality]\label{lem:duality}
$\Delta(A)$ is the annihilator of $I_2$ under this pairing.
\end{lemma}

\begin{proof}
Finite-dimensional fact: if $X$ and $Y$ are subspaces of a finite-dimensional space $V$, the restrictions
to $Y$ of the functionals vanishing on $X$ are exactly the functionals on $Y$ vanishing on $X\cap Y$. (A
functional $g$ on $Y$ vanishing on $X\cap Y$ defines $f(x+y)=g(y)$ on $X+Y$, well defined, and $f$ extends
to $V$.) Apply it with $V=(\mathfrak m/\mathfrak m^3)^s$, $X=\operatorname{Im}L_2$ and
$Y=(\mathfrak m^2/\mathfrak m^3)^s$. Via $T$, $V$ is the space of tuples of polynomials in $e$ with zero
constant term modulo $(e)^3$, and a functional on $V$ is a pair $(\nu,\Lambda)$, $\nu_i\in F^k$ pairing with
the linear coefficients and $\Lambda$ as above; its restriction to $Y$ is $\Lambda$. Write an element of
$(\mathfrak m/\mathfrak m^3)^M$ as $\sum_je_jP_j+\sum_{j\le l}e_je_lQ_{jl}$ with $P_j,Q_{jl}\in\Mat_{d+1}(F)$.
Since $T(p_{i\alpha}\otimes p_{i\beta})=p_{i\alpha}\sigma(p_{i\beta})$ and
$\sigma(p_i)\equiv p_i+\sum_le_lD_lp_i$ modulo $(e)^2$, the $i$-th entry of its image under $L_2$ is
\[
 \sum_je_j\,p_i^TP_jp_i+\sum_{j,l}e_je_l\,p_i^TP_jD_lp_i+\sum_{j\le l}e_je_l\,p_i^TQ_{jl}p_i .
\]
Write $\langle Z,W\rangle=\operatorname{tr}(Z^TW)$ for matrices, so that $p^TZq=\langle Z,pq^T\rangle$. The
pairing of $(\nu,\Lambda)$ with this image is
$\sum_j\bigl\langle P_j,\sum_i(\nu_i)_jp_ip_i^T+\sum_i\sum_l(\Lambda_i)_{jl}p_i(D_lp_i)^T\bigr\rangle
+\sum_{j\le l}\bigl\langle Q_{jl},\sum_i(\Lambda_i)_{jl}p_ip_i^T\bigr\rangle$ (for $j<l$ the form
$c\,e_je_l$ has matrix entries $c/2$ at $(j,l)$ and $(l,j)$, and $\Lambda_i$ is symmetric). So
$(\nu,\Lambda)$ vanishes on $\operatorname{Im}L_2$ if and only if $\Lambda$ satisfies (D1) and
$X_j=-\sum_i(\nu_i)_jp_ip_i^T$ for every $j$. The $\Lambda$ for which some $\nu$ exists are exactly those
satisfying (D1) and (D2).
\end{proof}

\begin{corollary}[Dual form of the 2-jet test, at every quadric rank]\label{cor:jetdual}
If $A$ affinely spans $\R^d$, the field criterion over $B/\mathfrak m^3$ (the 2-jet test) is solvable if
and only if $\sum_i\operatorname{tr}(\Lambda_iG_i)=0$ for every $\Lambda\in\Delta(A)$. The ``if'' direction
needs no spanning hypothesis.
\end{corollary}

\begin{proof}
By Proposition~\ref{prop:finitetype}(iii) the test is solvable if and only if
$b\in\Phi(\mathfrak mB^M)+\mathfrak m^3B^s$, that is, if and only if the class $\bar b$ of $b$ lies in
$\operatorname{Im}L_2$ (lift a preimage to $\mathfrak m^M$). By Lemma~\ref{lem:obstruction},
$\bar b\in(\mathfrak m^2/\mathfrak m^3)^s$, so this says $\bar b\in I_2=\Delta(A)^\perp$
(Lemma~\ref{lem:duality} and finite dimension), and $\langle\Lambda,\bar b\rangle=-\tfrac12\sum_i
\operatorname{tr}(\Lambda_iG_i)$. The step from $b\in\Phi(\mathfrak mB^M)+\mathfrak m^3B^s$ to a solution
($S\otimes1-Q$) does not use the span.
\end{proof}

\begin{remark}[Corank one]\label{rem:corankjet}
Let $\mathrm{Rel}=F\lambda$. For a derivation $\partial=\sum_l\xi_lD_l$ we have $\partial p_i=(0,\partial a_i)$,
so $\Pi_\lambda(\xi)=0$ says exactly that $\sum_i\lambda_i\partial a_i=0$ and
$\sum_i\lambda_i(\partial a_i)a_i^T$ is symmetric: the linear conditions of Lemma~\ref{lem:A1}. By
Lemma~\ref{lem:Pi}(b), $\Delta(A)=\mathrm{Sym}^2(K)$ with $K=\ker\Pi_\lambda$, which is spanned by the
$\xi\xi^T$, $\xi\in K$, and $\sum_i\operatorname{tr}(\xi\xi^T\lambda_iG_i)=\sum_i\lambda_i|\partial a_i|^2$, the quantity $E$ of Lemma~\ref{lem:A1}.
A linear functional on $\mathrm{Sym}^2(K)$ that vanishes on every $\xi\xi^T$ is zero, so at quadric corank
one the pairing of Corollary~\ref{cor:jetdual} is nonzero on $\Delta(A)$ (for a spanning set: the 2-jet test
fails) if and only if some derivation of $F$ satisfies all the hypotheses of Lemma~\ref{lem:A1} with this
$\lambda$; then $A$ is not Ramsey by Lemma~\ref{lem:A1} directly (extend the derivation to $\R$), with no
spanning hypothesis. This is the corank-one statement used in \S\ref{sec:jet}.
\end{remark}

\begin{theorem}[First-order rigidity; ledger SR-C34]\label{thm:rigidity}
Let $A$ be as in the setting of this section (the unit sphere centred at $0$; no spanning hypothesis). If
$\Delta(A)=0$, then $A$ is Ramsey.

\smallskip\noindent\emph{Status:} \st{HAND}: written proof in this section (added in v1.1.1),
independently reviewed twice by non-author AI agents (the first review reconstructed the argument; a
second, independent review found every step sound). Scope \texttt{ANY}.
\end{theorem}

\begin{proof}
Take $U,V,\Psi,\beta$ as in Lemma~\ref{lem:smith}. By Lemma~\ref{lem:local}(d), $\mathcal O/\mathfrak m^3\mathcal O=
B/\mathfrak m^3$, so $U$ and $V$ reduce to invertible matrices $\bar U,\bar V$ over $B/\mathfrak m^3$, and
$L_2$ is the restriction of $\Phi$ modulo $\mathfrak m^3$ to $(\mathfrak m/\mathfrak m^3)^M$, which $\bar V$
preserves. Hence
\[
 \operatorname{Im}L_2=\bar U\bigl((\mathfrak m/\mathfrak m^3)^\rho\oplus W\bigr),\qquad
 W:=\bigl(\mathfrak m\operatorname{Im}\Psi+\mathfrak m^3\mathcal O^\kappa\bigr)/\mathfrak m^3\mathcal O^\kappa
 \subset(\mathfrak m^2/\mathfrak m^3)^\kappa,
\]
the inclusion because $\Psi$ has entries in $\mathfrak m$. $\bar U$ is invertible and preserves
$(\mathfrak m^2/\mathfrak m^3)^s$, so
$I_2=\bar U\bigl((\mathfrak m^2/\mathfrak m^3)^\rho\oplus W\bigr)$ and, by Lemma~\ref{lem:duality},
\[
 \dim_F\Delta(A)=s\,\tfrac{k(k+1)}2-\dim I_2=\kappa\,\tfrac{k(k+1)}2-\dim_FW .
\]
If $\Delta(A)=0$, then $W=(\mathfrak m^2/\mathfrak m^3)^\kappa$, that is,
$Z=\mathfrak m\operatorname{Im}\Psi+\mathfrak mZ$ for $Z:=\mathfrak m^2\mathcal O^\kappa$, a finitely generated
$\mathcal O$-module containing $\mathfrak m\operatorname{Im}\Psi$. Nakayama's lemma, applied to
$Z/\mathfrak m\operatorname{Im}\Psi$, gives $Z=\mathfrak m\operatorname{Im}\Psi$. By Lemma~\ref{lem:smith},
$\beta\in Z$, so $b\in\Phi(\mathfrak m\mathcal O^M)$, and $A$ is Ramsey by
Proposition~\ref{prop:finitetype}(i).
\end{proof}

\noindent The theorem does not say that $\Delta(A)\ne0$ implies non-Ramsey: by
Corollary~\ref{cor:jetdual} the 2-jet test can pass with $\Delta(A)\ne0$, and then
Conjecture~\ref{conj:C1} is the open question.

\begin{corollary}[Special cases]\label{cor:rigidcases}
\begin{enumerate}
\item[(a)] ($\kappa=0$.) If the matrices $p_ip_i^T$ are linearly independent over $F$, $A$ is Ramsey. This
is OpenAI's quadric-independence criterion \cite[Prop.~7.3]{OAI172} (Remark~\ref{rem:quadric}); in this
case $\Psi$ is empty, and their lifting lemma \texttt{lift\_\allowbreak{}kernel} is the corresponding
order-zero argument.
\item[(b)] ($k=0$.) If $F$ is algebraic, $\Delta(A)=0$ and $A$ is Ramsey: Theorem~\ref{thm:A}.
\item[(c)] (The circle, one parameter.) Let $d=2$, let $F$ have transcendence degree $1$, and let $A$ avoid
the pole $(-1,0)$ (a rational rotation arranges this without changing $F$). With $t$, $\delta=d/dt$,
$s_i$, $\delta_i$ and the $n\times8$ matrix $N$ as in Theorem~\ref{thm:rank} ($n=s$), $\Delta(A)$ is
isomorphic to the left kernel of $N$. So if $\operatorname{rank}_FN=n$, $A$ is Ramsey: this is
Theorem~\ref{thm:rank} without its hypotheses that $F$ is regular over $\Q$ and that $5\le n\le8$.
\end{enumerate}
\end{corollary}

\begin{proof}
(a) and (b) are Lemma~\ref{lem:Pi}(b) and the definition. (c) Here $k=1$, $D_1=\delta$ (a basis of
$\mathrm{Der}_\Q(F)$, which is one-dimensional), and by Lemma~\ref{lem:Pi}(b) $\Delta(A)$ is isomorphic to
$\{\lambda\in\mathrm{Rel}:\Pi_\lambda(1)=0\}$. With $w_i=1+s_i^2$ and $v(s)=(1+s^2,1-s^2,2s)$, $p_i=v(s_i)/w_i$.
Put $\lambda'_i=\lambda_i/w_i^2$, a bijective rescaling. The six entries of $v(s)v(s)^T$ span the
polynomials of degree at most $4$ in $s$, so $\lambda\in\mathrm{Rel}$ if and only if
$\sum_i\lambda'_is_i^m=0$ for $m=0,\dots,4$. Next, $p_i(\delta p_i)^T=v_i(\delta v_i)^T/w_i^2$ plus a
multiple of $v_iv_i^T$, and $\delta v_i=\delta_iv'(s_i)$ with $v'(s)=(2s,-2s,2)$; the entries of
$\operatorname{asym}(v(s)v'(s)^T)$ are $-2s$, $1-s^2$ and $1+s^2$, which span the polynomials of degree at
most $2$. So $\Pi_\lambda(1)=0$ if and only if $\sum_i\lambda'_i\delta_is_i^m=0$ for $m=0,1,2$. Together:
$\lambda\in\mathrm{Rel}$ and $\Pi_\lambda(1)=0$ if and only if $\lambda'N=0$.
\end{proof}

\begin{remark}[Larger fields]\label{rem:largerfield}
The theorem may be applied with any finitely generated subfield $F'\subset\R$ containing the coordinates in
place of $F$ (with $F'\otimes_\Q F'$, its diagonal ideal and a basis of $\mathrm{Der}_\Q(F')$). The proofs above
use only that the coordinates lie in the finitely generated field, and the proof of
Proposition~\ref{prop:finitetype}(i) gives $P$ over $F'\otimes_\Q F'$ with $\Phi(P)=0$ and $m(P)=cS$,
$c\in F'^\times$. Its image under $F'\otimes_\Q F'\to\R\otimes_\Q\R$, multiplied by $c^{-1}\otimes1$, satisfies
OpenAI's real criterion; their \texttt{realCriterion\_\allowbreak{}fieldCriterion} and
\texttt{fieldCriterion\_\allowbreak{}ramsey} then give that $A$ is Ramsey.
\end{remark}

\begin{remark}[What the theorem does not say]\label{rem:silent}
\begin{enumerate}
\item[(a)] It is not a characterization, and $\Delta(A)$ depends on the coordinates, not only on the
congruence class. Suppose a nonzero derivation $\partial=\sum_l\xi_lD_l$ satisfies
$\partial p_i=Xp_i+\mu_ip_i$ for all $i$, with one matrix $X$ over $F$ and scalars $\mu_i\in F$; infinitesimal
M\"obius motions are of this form. Then
$\Pi_\lambda(\xi)=\operatorname{asym}\bigl((\sum_i\lambda_ip_ip_i^T)X^T+\sum_i\lambda_i\mu_ip_ip_i^T\bigr)=0$
for every $\lambda\in\mathrm{Rel}$, so if $\kappa\ge1$ the tuple $C^1=\xi\xi^T$, $C^r=0$ for $r\ge2$, is a
nonzero element of $\Delta(A)$. For example, if $a_i=Rr_i$ with rational $r_i$, $d$ of them linearly independent, and a
rotation $R$ whose entries generate a transcendental field, every derivation is of this form, so the
theorem is silent whenever $\kappa\ge1$, although $A$ is congruent to a rational set and is Ramsey by
Theorem~\ref{thm:A}. The theorem may be applied to any congruent copy, and by Theorem~\ref{thm:C} to any
M\"obius image.
\item[(b)] Not proved here (ledger SR-C34, consistent on sample only): the one-parameter threshold law
(configurations in general position on $S^{d-1}$ with one parameter are Ramsey if and only if they have
fewer than $(d+1)^2$ points), and the claim that the Ramsey sets with $\Delta=0$ on $S^2$ in the ledger's
sample are not two-group.
\end{enumerate}
\end{remark}

\begin{proposition}[Corank one; ledger SR-C33(b),(c)]\label{prop:corankone}
Suppose $\kappa=1$, $\mathrm{Rel}=F\lambda$. Choose $s-1$ columns $J$ of $\Phi$ whose reduction has rank
$s-1$, let $\beta_J=\det[\Phi_J\mid b]\in B$, and let $I\subset\mathcal O$ be the ideal generated by the
$s\times s$ minors of $\Phi$.
\begin{enumerate}
\item[(a)] $b\in\Phi(\mathfrak m\mathcal O^M)$ if and only if $\beta_J\in\mathfrak mI$ (membership in
$\mathcal O$). So $A$ is Ramsey if $\beta_J\in\mathfrak mI$, and, if $A$ affinely spans $\R^d$, only then.
$I=0$ is allowed, and then the condition is $\beta_J=0$ in $\mathcal O$ (the map $B\to\mathcal O$ need not
be injective when $B$ is not a domain).
\item[(b)] If $k=1$ (so $\mathcal O$ is a discrete valuation ring, with order function
$\operatorname{ord}$) and $I\ne0$, the condition is $\operatorname{ord}\beta_J>\operatorname{ord}I$.
\item[(c)] If $k=1$, let $D=D_1$ and $E=\sum_i\lambda_i|Da_i|^2$. Then $\operatorname{ord}\beta_J\ge2$, with
equality if and only if $E\ne0$; and $I\subset\mathfrak m^2\mathcal O$ if and only if $D$ satisfies the linear
conditions of Lemma~\ref{lem:A1} ($\sum_i\lambda_iDa_i=0$ and $\sum_i\lambda_i(Da_i)a_i^T$ symmetric).
Consequently: if the linear conditions fail, $\Delta(A)=0$ and $A$ is Ramsey; if they hold and $E\ne0$, $A$
is not Ramsey (Lemma~\ref{lem:A1}); if they hold and $E=0$ (the \emph{doubly degenerate} case) the 2-jet test
passes, and a non-Ramsey example requires $I\subset\mathfrak m^3\mathcal O$.
\end{enumerate}
\end{proposition}

\begin{proof}
(a) Take $U,V,\Psi$ as in Lemma~\ref{lem:smith}; $\Psi=\psi$ is a single row. By the Cauchy--Binet formula
the ideal of $s\times s$ minors is unchanged by invertible row and column operations, and the $s\times s$
minors of $\begin{psmallmatrix}I_{s-1}&0\\0&\psi\end{psmallmatrix}$ are $0$ and the entries of $\psi$. So
$I=\operatorname{Im}\psi$, the ideal generated by the entries of $\psi$. Next,
$\det(U)^{-1}\beta_J=\det[U^{-1}\Phi_J\mid U^{-1}b]$. The columns $J$ of $U^{-1}\Phi$ are
$\begin{psmallmatrix}V'\\\psi V''\end{psmallmatrix}$, where $V'$ and $V''$ are the first $s-1$ and the
remaining rows of the columns $J$ of $V$; the reduction of this block has rank $s-1$ and its last row
reduces to $0$, so $\det V'$ is a unit. Write $U^{-1}b=(x,\beta)$ with $x\in\mathfrak m^2\mathcal O^{s-1}$.
By the Schur complement,
\[
 \det(U)^{-1}\beta_J=\det(V')\bigl(\beta-\psi V''V'^{-1}x\bigr),
\]
and $\psi V''V'^{-1}x\in\mathfrak m^2I\subset\mathfrak mI$ (and $\subset\mathfrak m^3\mathcal O$). So
$\beta_J\in\mathfrak mI$ if and only if $\beta\in\mathfrak mI=\mathfrak m\operatorname{Im}\Psi$, which is
Lemma~\ref{lem:smith}; the Ramsey statements are Proposition~\ref{prop:finitetype}(i),(ii).

(b) If $I=\pi^c\mathcal O$, then $\mathfrak mI=\pi^{c+1}\mathcal O$.

(c) By the display, $\beta_J\equiv u\beta$ modulo $\mathfrak m^3\mathcal O$ with $u$ a unit. By
Lemma~\ref{lem:smith} the last row of $U_0^{-1}$ is $c_0\lambda$ with $c_0\in F^\times$, so
$\beta\equiv c_0\sum_i\lambda_ib_i$ modulo $\mathfrak m^3\mathcal O$, and by Lemma~\ref{lem:obstruction} and
Lemma~\ref{lem:local}(c), $\sum_i\lambda_ib_i\equiv-\tfrac12E\,\pi^2$. So $\operatorname{ord}\beta_J\ge2$,
with equality if and only if $E\ne0$. For $I$: in a discrete valuation ring, $I\subset\mathfrak m^2$ if and only
if $\mathfrak mI\subset\mathfrak m^3$, that is, $W=0$ in the proof of Theorem~\ref{thm:rigidity}; by the
dimension formula there ($\kappa=k=1$) this holds if and only if $\Delta(A)\ne0$, that is (Lemma~\ref{lem:Pi}),
$\Pi_\lambda(1)=0$, which is the pair of linear conditions (Remark~\ref{rem:corankjet}). The consequences:
if the conditions fail, $\Delta(A)=0$; if they hold and $E\ne0$, Lemma~\ref{lem:A1} applies with $D$
extended to $\R$; if they hold and $E=0$, the pairing of Corollary~\ref{cor:jetdual} vanishes on
$\Delta(A)=F$, and if $\operatorname{ord}I=2$ then $\operatorname{ord}\beta_J\ge3>\operatorname{ord}I$ and $A$
is Ramsey by (b).
\end{proof}

\noindent\emph{What remains open here} (ledger SR-C33(d)): whether doubly degenerate corank-one
configurations with $I\subset\mathfrak m^3\mathcal O$ and $\operatorname{ord}\beta_J\le\operatorname{ord}I$
exist; such a configuration would be a non-Ramsey set with no scalar certificate, and a counterexample to
Conjecture~\ref{conj:C1}.

\section{Conjectures}\label{sec:conj}

\begin{conjecture}[C1, second-order completeness; ledger SR-C1]\label{conj:C1}
A finite spherical set is Ramsey iff its 2-jet system is solvable with $c=1$. Evidence:
\S\ref{sec:jet}, \st{PROBE}.
\end{conjecture}

\begin{conjecture}[C3, generic thresholds; ledger SR-C3]\label{conj:C3}
A generic $n$-subset of $S^{d-1}$ is non-Ramsey iff $n\ge N(d)$, with $N(2)=7$ and $N(d)=(d+1)(d+2)/2$ for
$d\ge3$. Evidence \st{PROBE}: circle $6/7$, $S^2$ $9/10$; first failures at $15$ ($S^3$) and $21$ ($S^4$)
with four or five free points. On the failing side the $S^2$ boundary cells are certified (Theorem~\ref{thm:NR}, \st{CERT}),
and the circle boundary cell is proved by hand (Proposition~\ref{prop:PD}, \st{HAND}, reviewed (AI)). Related (C4): algebraic points plus at most three
arbitrary points are Ramsey on every sphere, and four can fail. The first half follows from
Theorem~\ref{thm:C2} (Remark~\ref{rem:C4}); four free points fail on $S^1$, $S^2$, $S^3$
(Remark~\ref{rem:C4}), and \st{PROBE} for $5\le d\le9$. Ledger entry SR-C17 merges C3 and C4 into one
phase-diagram conjecture for $m$ rational points in general position plus $j$ independent points on
$S^{d-1}$, $d\ge3$: non-Ramsey iff $j\ge4$ and $m+j\ge(d+1)(d+2)/2$.

\emph{Status in version~1.1: proved, mostly by prior work} (ledger SR-C3; this is bookkeeping, not a new
result). On the circle the two halves are Theorem~\ref{thm:Ba} and Proposition~\ref{prop:PD}. For
$d\ge3$ the Ramsey half ($n<N(d)$) is a corollary of OpenAI's Proposition~7.3
(Remark~\ref{rem:quadric}), and the non-Ramsey half is P\'alv\"olgyi's Theorem~A.12 \cite{Palvolgyi},
from an appendix that its author labels unchecked (its proof was read by our AI reviewer), and
independently the case $m=0$ of ledger entry SR-C30 (\st{SKETCH}, reviewed once by an AI agent). For
SR-C17 the passing side is proved (Remark~\ref{rem:C4} for $j\le3$, Proposition~7.3 for $m+j\le D_d$);
the failing side has exact certificates for specific rational points for $d\le7$ and a sketch in every
dimension (SR-C30).
\end{conjecture}

\begin{conjecture}[R1, rigidity; ledger SR-C6]\label{conj:R1}
For $n\ge9$ points on a circle, every positive-dimensional family of Ramsey configurations lies in a
two-M\"obius family (Theorem~\ref{thm:C2}, $d=2$). Theorem~\ref{thm:rigid} proves this for families
with regular coordinate field that contain a five-point group with rational cross-ratios, and
Conjecture~\ref{conj:count}(ii) predicts it for rational families. The bound $n\ge9$ is needed:
Theorem~\ref{thm:NTG} and Corollary~\ref{cor:radian8} give non-two-group families with $7$ and $8$ points.
\end{conjecture}

\begin{conjecture}[Expected-dimension law on the circle]\label{conj:count}
Let $A\subset S^1$, $|A|=n$, have a purely transcendental coordinate field, and let $k$ be the transcendence
degree of the field generated by its four-point cross-ratios (a M\"obius invariant).
\begin{enumerate}
\item[(i$'$)] If $(n-5)(k+1)\le6$ and $A$ is non-degenerate, then $A$ is Ramsey. For $k=1$ non-degenerate
means $\operatorname{rank}N=n$, and (i$'$) is Theorem~\ref{thm:rank}; for $k=2$ it means that the 2-jet system is
solvable, and (i$'$) is then an instance of Conjecture~\ref{conj:C1}.
\item[(ii)] If $(n-5)(k+1)>6$, then $A$ is Ramsey if and only if it is two-group.
\end{enumerate}
For $n=7$ this would say: Ramsey iff two-group or ($k\le2$ and non-degenerate). Evidence: the count is that of
the reduced 2-jet system ($n\,k(k+1)/2$ equations against $3k+5\,k(k+1)/2$ unknowns), and is derived, not
fitted, for $k=1$. (ii): $0$ mismatches on $263$ fresh seeded configurations \st{PROBE}, and $176$
one-parameter $9$-point sets that are not two-group have an exactly nonzero $9\times9$ evaluation determinant,
hence are not Ramsey (\st{CERT} per instance). The first form of (i), without the non-degeneracy hypothesis,
is refuted at probe level: the mirror-symmetric rational family $\{0,\pm a,\pm b,\pm c\}$ has $k=2$ and fails
the 2-jet test (\st{PROBE}, via Proposition~\ref{prop:J2}). Square-root fields are excluded: P\'alv\"olgyi's
heptagon has $n=7$, $k=1$ and is not Ramsey. Ledger entry SR-C23.
\end{conjecture}

\begin{conjecture}[T1; ledger SR-C9]\label{conj:T1}
Ramsey-ness of explicit configurations can be equivalent to open transcendence problems
(Corollary~\ref{cor:radian}); conversely, Schanuel's conjecture would make Ramsey-ness decidable for
exponential-algebraic configurations.
\end{conjecture}

\begin{conjecture}[M1, measurable Graham; ledger SR-C5]\label{conj:M1}
Let $A$ be a finite spherical set.
\begin{enumerate}
\item[(M1)] For every $r$ there is $D$ such that every Lebesgue-measurable $r$-colouring of $\R^D$ is
constant on some congruent copy of $A$.
\item[(M1w)] Moreover the monochromatic copies have positive measure: for every $r$ there is $D$ such that,
for every Lebesgue-measurable $r$-colouring $\chi$ of $\R^D$, the set of $(v,g)\in\R^D\times\Ort(D)$ for
which $\chi$ is constant on $v+gA$ has positive measure (Lebesgue times Haar).
\item[(M1cat)] (Category twin.) For every $r$ there is $D$ such that whenever $U_1,\dots,U_r$ are disjoint
open subsets of $\R^D$ with dense union, some $U_j$ contains a congruent copy of $A$.
\end{enumerate}
\end{conjecture}

\noindent\st{SKETCH} (not reviewed). M1w is unchanged by null modifications of $\chi$, so it
reduces to Borel colourings, and it is a $\Pi^1_2$ statement; if ZFC proves it, Shoenfield absoluteness
gives Graham's conjecture in Solovay's model, where every set is measurable, although the conjecture fails
in ZFC \cite{Palvolgyi}. (Graham in Solovay's model does not obviously imply M1w.) M1cat would give
Graham's conjecture in Shelah's model of ZF$+$DC$+$BP, with no inaccessible cardinal. A counterexample to
M1w would be some $r$ and, for every $D$, a Borel $r$-colouring of $\R^D$ whose monochromatic copies of
$E_1$ (Corollary~\ref{cor:radian}) form a null set. \cite{OAI172} explicitly assumes no measurability. We
found no statement of M1 in our searches. Zapletal \cite{Zapletal} works in choiceless models
(ZF$+$DC) with countable chromatic numbers of triangle hypergraphs and says nothing about finite
colourings, Graham's conjecture or M1. The nearest literature: measurable chromatic numbers \cite{Falconer},
\cite{ShelahSoifer,Payne}; density theorems for simplices and related configurations
\cite{Bourgain,CastroSilva} and on spheres \cite{GuruswamiLi}; density Graham over $\mathbb F_p$ for all
spherical configurations \cite{Sun}; P\'alv\"olgyi's question about a density version
\cite[\S7]{Palvolgyi-ncs}; and \cite{OAI158}, where measurability does not matter. Current Euclidean density
theory covers simplices, their products and inductive configurations, so we see no route to M1 now.

\begin{remark}[Three lemmas around M1]\label{rem:M1lemmas}
\begin{enumerate}
\item[(L1)] \st{HAND} A Lebesgue-measurable $\Q$-quadratic map $\R\to\R$ is $x\mapsto cx^2$, and a
Lebesgue-measurable $\Q$-derivation of $\R$ is $0$: measurable additive maps are linear, and $D(x^2)=2xD(x)$
then forces $c=0$. \st{SKETCH} So the separating data in OpenAI's necessity argument \cite[\S2]{OAI172}
cannot be measurable for a spherical set; restricted to measurable data, that mechanism detects only
non-sphericity.
\item[(L3)] \st{HAND} (via Theorem~\ref{thm:A}). ``Robust Graham'' holds in ZFC: for every finite spherical
$A$, every $\varepsilon>0$ and every $r$ there is $D$ such that every $r$-colouring of $\R^D$ is constant on a
congruent copy of some spherical $A'$ within $\varepsilon$ of $A$, because spherical sets with algebraic
coordinates are dense. So a counterexample to M1 cannot be robust under perturbation.
\item[(L4)] \st{HAND} (compactness; not reviewed). Let $r_0$ be a number of colours for which $E_1$ is not
Ramsey, and let $\theta_n\to1$ with $e^{i\theta_n}$ algebraic, so that $E_{\theta_n}$ is Ramsey. Let $N_n$ be
the least size of a finite set $W$ in some $\R^D$ such that every $r_0$-colouring of $W$ is constant on a
congruent copy of $E_{\theta_n}$. Then $N_n\to\infty$. Otherwise, along a subsequence the witnesses have a
fixed size, a fixed family of index tuples forming copies, and (minimal witnesses are connected, so after
translation they are bounded) convergent points; the limit set is a finite witness for $E_1$, because
limits of copies of $E_{\theta_n}$ are copies of $E_1$ and a colouring of the limit pulls back. So any proof of
M1 for $E_1$ cannot pass through bounded finite witnesses for nearby Ramsey sets.
\end{enumerate}
\end{remark}

\begin{remark}[Zero--one law]
\st{SKETCH}, scope \texttt{GENERIC\_OF}. In an irreducible family of spherical configurations, generic
members share one verdict (their coordinate fields are isomorphic compatibly with the coordinates), and
Ramsey members are dense (Theorem~\ref{thm:A}).
\end{remark}

\begin{remark}[Definable and polynomial-factor colourings]\label{rem:definable}
These statements test candidate counterexamples to Conjecture~\ref{conj:M1} and do not approach it.
\begin{enumerate}
\item[(a)] \st{HAND}, reviewed (AI); o-minimal folklore \cite{vdDries}, recorded here as an observation, not
as a result. \emph{Big-ball lemma.} Let $\mathcal R$ be an o-minimal expansion of the
real field (semialgebraic sets with arbitrary real parameters, $\R_{\exp}$, $\R_{\mathrm{an}}$, \dots), $D\ge2$,
and $\R^D=C_1\sqcup\dots\sqcup C_r$ with every $C_j$ definable in $\mathcal R$. Then some single $C_j$ contains
balls of every radius, hence a congruent copy of every bounded subset of $\R^D$, spherical or not.
\emph{Proof.} The union $\mathrm{Bd}$ of the boundaries $\partial C_j$ is definable of dimension $\le D-1$
\cite[Ch.~4]{vdDries}. Under $\varphi(x)=x/\sqrt{1+|x|^2}$, which maps $\R^D$ onto the open unit ball, the
points of $\overline{\varphi(\mathrm{Bd})}$ on the unit sphere lie in the frontier of $\varphi(\mathrm{Bd})$,
of dimension $\le D-2$, so some $\omega$ on the unit sphere has a neighbourhood missing
$\overline{\varphi(\mathrm{Bd})}$. Then a truncated cone $\{x:|x|>R_0,\ |x/|x|-\omega|<\epsilon\}$ misses
$\mathrm{Bd}$; it is connected ($D\ge2$), so it lies in one class, and it contains balls of every radius.
For semialgebraic classes there is an elementary variant: $\mathrm{Bd}\subset Z(g)$ for a nonzero polynomial
$g$, and if the top-degree form of $g$ does not vanish at $\omega$, a truncated cone around $\omega$ misses
$Z(g)$. So for definable colourings the Ramsey property holds trivially for every bounded
configuration, and this class cannot tell spherical sets from non-spherical ones.
\item[(b)] \st{HAND} modulo two cited theorems, reviewed (AI). \emph{Polynomial factors.} For all $d,m,k$ there
is $N(d,m,k)$ such that for $D\ge N$, any $P_1,\dots,P_m\in\R[x_1,\dots,x_D]$ of degree $\le d$ are
simultaneously radial (functions of $|x|$) on some $k$-dimensional linear subspace $V$. Hence every colouring
$f\circ(P_1,\dots,P_m)$, with $f:\R^m\to[r]$ arbitrary (not even measurable) and any $r$, is constant on every
sphere of $V$ centred at $0$, and contains a monochromatic congruent copy of every finite spherical set of
affine dimension $\le k$. \emph{Proof.} Apply Dol'nikov and Karasev's Theorem~1 (an even-degree form is
proportional to $|x|^{\deg}$ on a subspace of prescribed dimension, given enough variables) and their
Theorem~5 (an odd polynomial map to $\R^m$ vanishes on such a subspace) \cite{DolnikovKarasev} to the
homogeneous components one at a time, inside a decreasing chain of subspaces; both conclusions pass to
subspaces. The cited theorems are read at the level of their statements; $N$ is not effective. The same
holds, with a bound $N(d,s,m,k)$, for $m$ bracket-polynomial expressions built with at most $s$ gates
$+,\times,\lfloor\cdot\rfloor$ over polynomial leaves of degree $\le d$, by induction on the floor depth; it
does not cover families whose number of brackets grows with $D$. The ingredients are Dol'nikov and
Karasev's \cite{DolnikovKarasev}; Bukh and Karasev \cite[\S2]{BukhKarasev} already link
Dvoretzky- and Knaster-type theorems to Euclidean Ramsey problems, and Gromov's ``Hilbert
symmetrization'' and ``Ramsey--Dvoretzky'' discussion \cite[\S1.12]{Gromov} is prior art for this
radialisation idea. Our contribution is only the observation that it disposes of this class of
candidate colourings.
\item[(c)] \st{HAND}, reviewed (AI). \emph{Cube lemma.} For $D\ge3$ the vectors with coordinates
$u_i=\sqrt{2/D}\cos(2\pi i/D)$, $v_i=\sqrt{2/D}\sin(2\pi i/D)$ are orthonormal, so every coordinate of
$c+u\cos\phi+v\sin\phi$ stays within $\sqrt{2/D}$ of $c_i$, and a unit circle fits in a closed cube of side
$2\sqrt{2/D}$. Hence a colouring with a class containing such a cube has a monochromatic copy of every finite
subset of a unit circle, and a colouring constant on the (half-open) cells of a cubical grid of side $h$ is
defeated once $D>8/h^2$.
\item[(d)] \textbf{Open} (ledger SR-C28). Graham's conjecture for \emph{locally semialgebraic} colourings,
whose classes meet every ball in a semialgebraic set. This class contains the colourings
$x\mapsto\lfloor c|x|^2\rfloor\bmod r$ of \cite{EGMRSS}, which defeat every non-spherical set, so it does see
sphericity; Conjecture~\ref{conj:M1} implies it. By (b) and (c), a counterexample family needs factor
complexity growing with $D$ and cells finer than $D^{-1/2}$. \st{SKETCH}: in density form the question is
equivalent to the density form of M1 for Borel colourings.
\end{enumerate}
A literature search found no statement of Graham's conjecture for definable colourings
(ledger SR-C25 to SR-C28).
\end{remark}

\section{Tailings}\label{sec:tailings}

Killed ideas and harness lessons, recorded so that others need not repeat them.
\begin{itemize}
\item Poncelet chains are not special: they fail at seven points like random three-parameter families
\st{PROBE}. ``Two groups only on Hopf spheres'' is refuted by $S^3$, and now by
Theorem~\ref{thm:C2}. A naive count predicting that the free-point threshold grows with $d$ is refuted (it
is $3$ for $d=2,\dots,9$, \st{PROBE}).
\item ``Matrix weights are strictly stronger than scalar ones'': withdrawn for its original witness
(\S\ref{sec:jet}). For certificates it is true, with new witnesses (ledger SR-C32); for colourings it is
not known, since a matrix-weighted Lemma~\ref{lem:A1} colours those witnesses (ledger SR-C35, a sketch).
\item ``A $7$-point cyclic set is Ramsey iff it is two-group'': refuted (Theorem~\ref{thm:NTG}). Its
refinement to cross-ratio transcendence degree $\ge2$ is refuted at probe level (the threshold is $3$;
Conjecture~\ref{conj:count}). ``One-parameter cyclic sets of at most eight points are Ramsey'' is false as
stated: P\'alv\"olgyi's heptagon is a one-parameter set; the rank condition of Theorem~\ref{thm:rank} is
what fails there.
\item A predicted Ramsey generic member of the mirror family $\{0,\pm a,\pm b,\pm c\}$, which would have
shown that Ramsey-ness is not closed under specialization on the circle, failed the 2-jet test
\st{PROBE}; off the circle the question is settled by OpenAI's Proposition~7.3
(Remark~\ref{rem:quadric}; ledger SR-C16).
\item The generic-entry quaternionic recipe for M\"obius graph forms does not extend to $\Cl_{0,4}$
(\S\ref{sec:lowdim}); the Lipschitz normal form is needed.
\item Three probe traps: integer division producing floats, and a point at the stereographic pole, each
produced false passes before being fixed; a retry routine whose ``no verdict'' return was counted as
agreement inflated an agreement count (corrected; no pass/fail flip was found).
\item A random sampler of one-parameter configurations with square roots produced non-real configurations
($507$ of $1557$); all evidence above is restricted to real ones, decided exactly.
\end{itemize}

\section{AI use and review disclosure}\label{sec:disclosure}

The author, Will Strinz, is the only human contributor. He directed the work and
made its decisions; he is not an expert in Euclidean Ramsey theory. The statements, proofs and probes
first came out of an exploratory conversation with Claude Opus~5.5 (Anthropic, in claude.ai) on
2026-10-11/12. The campaign was then coordinated by Claude Opus~5.5 in Claude Code, which dispatched
Claude Opus~5.5 agents (the Lean proofs and ropes, the exact certificates and probes including the
Clifford construction of \S\ref{sec:twogroups}, the reviews below, and this note) and Claude Sonnet~5.5
agents (the literature searches, hygiene checks of the probes, and three fresh reproductions). Every
machine check was run by these agents; the author did not re-run them.

All reviews are AI reviews, each by a fresh agent that had not built what it reviewed: a cold review of the
second-order test (Proposition~\ref{prop:J2}); adversarial reviews of the certificate glue for $S^1$--$S^4$
and the circle phase diagram, for $S^5$ and every $d$ (Theorem~\ref{thm:C2}), and for the $S^2$/$S^3$
certificates with the colouring lemma (Theorem~\ref{thm:NR}, Lemma~\ref{lem:A1}); non-builder fidelity
reviews of ropes AV1, AV3 and AV5 and of the Theorem~C package; fresh reproductions of rope AV1, of the
Theorem~C package, and of all five ropes of this release with its certificates (fast and slow sets)
and probes; a non-author review of Theorems~\ref{thm:rank} and \ref{thm:rigid}; and a skeptical
review of the statements of Remark~\ref{rem:definable}. The certificates of
Theorem~\ref{thm:NTG} were re-verified by a non-builder with independent code, and their glue was
reviewed (AI), as was the hypothesis check for P\'alv\"olgyi's Theorems~A.3 and A.7 in
Proposition~\ref{prop:PD}. The author read
the hand proof of Theorem~\ref{thm:C} line by line, and reviewed this note and the ledger before
release (not line by line); both are an owner's sign-off, not expert review. The corrections of
version~1.1 were found and written by Claude Opus~5.5 agents and reviewed by a separate AI agent
(including the genericity step of Remark~\ref{rem:quadric}), and the two additions listed under
``Changes in version~1.1'' were each reviewed by one non-author AI agent. For first-order rigidity that
reviewer reconstructed a complete proof, which version~1.1.0 summarized without printing it. In
version~1.1.1 the written proofs of \S\ref{sec:rigidity} were added by a Claude Opus~5.5 agent, following
that reconstruction and making each step explicit; a second, independent review by a Claude Opus~5.5 agent
that had not built them, and did not see the earlier review before reaching its verdict, found every step
sound, and its exact search for counterexamples over about fifteen families of configurations found no
non-Ramsey set with $\Delta(A)=0$. \textbf{No human expert has reviewed any part of this note.} Machine checks:
Lean and Comparator (Theorems~\ref{thm:A}, \ref{thm:Ba}, \ref{thm:Bb}, \ref{thm:C} with
Corollary~\ref{cor:mobius}(c),(d), and \ref{thm:D}); exact sympy identities and exact ranks (\st{CERT}
items, and the rank computation in Corollary~\ref{cor:radian8}); random
specializations modulo primes (\st{PROBE} items only). This follows the spirit of the disclosures in
\cite{Palvolgyi}, \cite{MR} and the vibemathed catalogue \cite{vibemathed}.

\appendix
\section{Verification}\label{app:lean}

\noindent\textbf{Ropes.} Each rope is a standalone Lean package (OpenAI's files vendored unmodified,
Apache-2.0, with NOTICE) checked by stock Comparator, under Landlock, against OpenAI's challenge
definitions, with negative controls (\texttt{sorry}, an extra axiom, a weaker statement, a redefined
constant, a weaker Mathlib-vocabulary statement). ``GR1'' is the release gate: a fresh-runner CI run, a
non-builder fidelity review, and a fresh-agent reproduction. The CI runs are this release's own
workflows on a private staging copy of the release repository, on a tree identical to the released
one; the release repository re-runs them on every push.

\medskip
{\small\noindent\begin{tabular}{@{}p{0.07\textwidth}p{0.1\textwidth}p{0.28\textwidth}p{0.22\textwidth}p{0.25\textwidth}@{}}
\toprule
Rope & Result & Lean theorems & Fresh-runner CI & GR1 \\
\midrule
AV1 & Thm~\ref{thm:A} & \texttt{ramsey\_of\_algebraic} & run \texttt{38493016206}: 5 accepted, 5 controls
refused & complete (v1.1 documentation fixes applied) \\
AV2 & Thm~\ref{thm:Ba} & \texttt{circle\_ramsey\_of\_le\_six} & run \texttt{38493016157}: all 10 as expected &
CI and fresh-agent reproduction passed; fidelity self-assessed, review open \\
AV3 & Thm~\ref{thm:D} & \texttt{radian\_configuration}, \texttt{radian\_not\_ramsey},
\texttt{radian\_ramsey\_iff}, \texttt{radian\_ramsey\_iff\_tan} & run \texttt{38493016161}: all 10 as expected &
complete (fidelity: faithful, documentation fixes applied; reproduction: reproduced, friction fixed) \\
AV4 & Thm~\ref{thm:C}, Cor.~\ref{cor:mobius} & \texttt{ramsey\_\allowbreak{}iff\_\allowbreak{}ramsey\_\allowbreak{}mobius\_\allowbreak{}full},
\texttt{ramsey\_\allowbreak{}mobius\_\allowbreak{}of\_\allowbreak{}algebraic}, \texttt{ramsey\_\allowbreak{}circ\_\allowbreak{}*} & run \texttt{38493016252}: all 12 as expected
& CI and fresh-agent reproduction passed for the rope; fidelity self-assessed; see the package \\
AV5 & Thm~\ref{thm:Bb} & \texttt{ramsey\_alg\_plus\_three} & run \texttt{38493016164}: all 10 as expected &
complete (fidelity: faithful, documentation fixes applied; reproduction: reproduced, friction fixed) \\
\cite{MR} & Thm~\ref{thm:C}, Cor.~\ref{cor:mobius} & the AV4 statements, ported to v4.35.0-rc2 & staging run
\texttt{38336666705}; released commit run \texttt{38339951762} & complete for the package (fidelity: faithful,
fixes applied; reproduction: reproduced, friction fixed) \\
\bottomrule
\end{tabular}}

\medskip\noindent All theorems build at the root with axioms \texttt{propext}, \texttt{Classical.choice},
\texttt{Quot.sound} only. The published package \cite{MR} checks with \texttt{lake comparator} and three
kernels (Lean, NanoDa, con-ron) and refuses five controls.

\noindent\textbf{Theorem C in Lean} (ASCII transliteration; the Unicode source is quoted verbatim in
\cite[\S3]{MR}; \texttt{Ramsey} and \texttt{Space} are OpenAI's):
\begin{verbatim}
theorem ramsey_iff_ramsey_mobius_full {s d : N} (a : Fin s -> Space d)
    (L : Matrix (Option (Fin d)) (Option (Fin d)) R)
    (hL : L^T * lorentzForm d * L = lorentzForm d) (ha : Function.Injective a)
    (hsph : forall i, a i in Metric.sphere (0 : Space d) 1) :
    Ramsey a <-> Ramsey (mobiusMap L o a)
\end{verbatim}

\noindent\textbf{Exact scripts} (Python 3.10, sympy 1.14; in \texttt{certificates/}, with expected logs):
\texttt{cert\_\allowbreak{}mobius.py} (Theorems~\ref{thm:Ba}, \ref{thm:Bb});
\texttt{cert\_\allowbreak{}s2.py} (about 2.3 min), \texttt{cert\_\allowbreak{}s2\_\allowbreak{}quat.py} (1--2 min), \texttt{g4\_\allowbreak{}test.py},
\texttt{cert\_\allowbreak{}clifford.py} (54--63 min), \texttt{cert\_\allowbreak{}clifford\_\allowbreak{}obs.py}, \texttt{cert\_\allowbreak{}clifford\_\allowbreak{}seed.py},
\texttt{cert\_\allowbreak{}clifford\_\allowbreak{}lorentz.py}, \texttt{cert\_\allowbreak{}clifford\_\allowbreak{}crosscheck.py} (corroboration; its coverage test is floating point) (Theorem~\ref{thm:C2});
\texttt{radian\_\allowbreak{}det.py} (Theorem~\ref{thm:D}); \texttt{dual\_\allowbreak{}cert.py} (J2-F7); \texttt{scalar\_\allowbreak{}exact.py}
(Theorem~\ref{thm:NR}); \texttt{cert\_\allowbreak{}nontwogroup.py} and \texttt{exact\_\allowbreak{}circle\_\allowbreak{}cert.py}
(Theorem~\ref{thm:NTG}); \texttt{rankN\_\allowbreak{}check.py} and \texttt{rankN\_\allowbreak{}sweep.py} (Theorem~\ref{thm:rank}, Corollary~\ref{cor:radian8});
\texttt{definable\_\allowbreak{}bigball.py}, \texttt{definable\_\allowbreak{}polyphase.py}, \texttt{definable\_\allowbreak{}crosscheck.py} and \texttt{exact\_\allowbreak{}intervals.py} (exact sanity checks for Remark~\ref{rem:definable});
\texttt{corank\_\allowbreak{}exact.py} (Remark~\ref{rem:quadric} and ledger entries SR-C29 to SR-C32),
\texttt{c31\_\allowbreak{}respecialize.py} and \texttt{c34\_\allowbreak{}det16.py} (ledger SR-C31, SR-C34; added in
version~1.1, with python-flint).
Modulo-$p$ probes, in \texttt{probes/}: \texttt{probe2.py}, \texttt{scalar\_\allowbreak{}cert.py}, the \texttt{e*.py}
experiments and the scripts behind the conjecture ledger (for example \texttt{c14\_\allowbreak{}rigid.py},
\texttt{c23\_\allowbreak{}count\_\allowbreak{}law.py}); the \texttt{validate.py} controls must pass before any
probe output is cited. They are \st{PROBE} evidence, not proof.

\noindent\textbf{Replay.} \texttt{REPRODUCE.md} of this release: each rope under \texttt{ropes/av1}--\texttt{av5}
is a standalone Lean package with its own \texttt{REPRODUCE} script (Comparator under Landlock on Linux); the
exact scripts run with \texttt{python -I certificates/<script>} and their output is compared with the
expected logs. The Theorem~C package: \texttt{REPRODUCE.md} of \cite{MR}. Repository and archive for this
note: \url{https://github.com/wstrinz/spherical-ramsey}, archived on Zenodo (the DOI is given in
the repository's \texttt{README.md} and \texttt{CITATION.cff}).

\begin{thebibliography}{99}
\bibitem{OAI172} OpenAI, \emph{A classification of finite Euclidean Ramsey configurations},
OpenAI Math Release preprint, September 23, 2026; github.com/openai/math, commit
\texttt{fd4aeeb2ee4fc729c18d98444fed42fd0529eeeb},
\texttt{preprints/A-classification-of-finite-Euclidean-Ramsey-configurations-September-23-2026/}.
\bibitem{OAI158} OpenAI, \emph{The Euclidean plane is not five-colorable}, OpenAI Math Release preprint, September 23,
2026 (release \#158); github.com/openai/math, commit \texttt{fd4aeeb2},
\texttt{preprints/The-Euclidean-plane-is-not-five-colorable-September-23-2026/}.
\bibitem{MR} W.~Strinz, \emph{mobius-ramsey: M\"obius invariance of the Ramsey property for finite
spherical sets (Lean~4)}, version 1.0.0, Zenodo, 2026, doi:10.5281/zenodo.23339523 (all versions:
doi:10.5281/zenodo.23339522); \url{https://github.com/wstrinz/mobius-ramsey}, release v1.0.0 at commit
\texttt{52284dd}. Includes the note \emph{Ramsey-ness of finite spherical sets is M\"obius invariant}.
\bibitem{Palvolgyi} D.~P\'alv\"olgyi, \emph{A cyclic non-Ramsey heptagon}, arXiv:2609.23327 (v1,
2026-09-20; no later version as of 2026-10-13).
\bibitem{Palvolgyi-ncs} D.~P\'alv\"olgyi, \emph{A nearcircumsphere-Ramsey theorem for solvable transitive
configurations}, arXiv:2608.10865.
\bibitem{JDJones} JD-Jones-ASES, \emph{six-concyclic-lean}, GitHub repository,
\url{https://github.com/JD-Jones-ASES/six-concyclic-lean}, root commit \texttt{0703d7d}, 2026-10-07.
\bibitem{Atimics} atimics, \emph{math2}, GitHub repository, \url{https://github.com/atimics/math2},
\texttt{EndToEnd/AlgebraicRamsey.lean}, commit \texttt{8ab51fa}, 2026-10-08.
\bibitem{Rogerhu} Rogerhu12, \emph{openai-math-extensions},
\texttt{papers/172-euclidean-ramsey/algebraic\_spherical\_configurations.tex}, commit \texttt{f8a2bb2},
2026-10-07, \url{https://github.com/Rogerhu12/openai-math-extensions}.
\bibitem{Eliott} eliottcassidy2000, \emph{math}, GitHub repository, \url{https://github.com/eliottcassidy2000/math},
audit note \texttt{04-computation/experiments/oai2\_20261006\_audits/audit\_A.md}, \S2 (theorem
THM-4559, conditional on \cite{OAI172}), commit \texttt{9cbd9ad}, 2026-10-07.
\bibitem{EGMRSS} P.~Erd\H{o}s, R.~L.~Graham, P.~Montgomery, B.~L.~Rothschild, J.~Spencer, E.~G.~Straus,
\emph{Euclidean Ramsey theorems I}, J. Combin. Theory Ser. A \textbf{14} (1973), no.~3, 341--363.
doi:10.1016/0097-3165(73)90011-3.
\bibitem{Graham94} R.~L.~Graham, \emph{Recent trends in Euclidean Ramsey theory}, Discrete Math. \textbf{136}
(1994), no.~1--3, 119--127. doi:10.1016/0012-365X(94)00110-5.
\bibitem{FR86} P.~Frankl, V.~R\"odl, \emph{All triangles are Ramsey}, Trans. Amer. Math. Soc.
\textbf{297} (1986), no.~2, 777--779. doi:10.1090/S0002-9947-1986-0854099-6.
\bibitem{FR90} P.~Frankl, V.~R\"odl, \emph{A partition property of simplices in Euclidean space}, J. Amer.
Math. Soc. \textbf{3} (1990), no.~1, 1--7. doi:10.1090/S0894-0347-1990-1020148-2.
\bibitem{Kriz91} I.~\Kriz, \emph{Permutation groups in Euclidean Ramsey theory}, Proc. Amer. Math. Soc.
\textbf{112} (1991), no.~3, 899--907. doi:10.1090/S0002-9939-1991-1065087-9.
\bibitem{Kriz92} I.~\Kriz, \emph{All trapezoids are Ramsey}, Discrete Math. \textbf{108} (1992), 59--62.
doi:10.1016/0012-365X(92)90660-8.
\bibitem{Cantwell} K.~Cantwell, \emph{All regular polytopes are Ramsey}, J. Combin. Theory Ser. A \textbf{114}
(2007), no.~3, 555--562. doi:10.1016/j.jcta.2006.08.001.
\bibitem{LRW12} I.~Leader, P.~A.~Russell, M.~Walters, \emph{Transitive sets in Euclidean Ramsey theory},
J. Combin. Theory Ser. A \textbf{119} (2012), no.~2, 382--396. doi:10.1016/j.jcta.2011.09.005.
\bibitem{LRW11} I.~Leader, P.~A.~Russell, M.~Walters, \emph{Transitive sets and cyclic quadrilaterals},
J. Comb. \textbf{2} (2011), no.~3, 457--462. doi:10.4310/JOC.2011.v2.n3.a6. arXiv:1012.5468.
\bibitem{GibbonsLuo} J.~C.~Gibbons, Y.~Luo, \emph{Colorings of the $n$-sphere and inversive geometry},
arXiv:1307.5351 (2013).
\bibitem{Ratcliffe} J.~G.~Ratcliffe, \emph{Foundations of Hyperbolic Manifolds}, 3rd ed., Graduate Texts in
Mathematics \textbf{149}, Springer, Cham, 2019. doi:10.1007/978-3-030-31597-9.
\bibitem{Beardon} A.~F.~Beardon, \emph{The Geometry of Discrete Groups}, Graduate Texts in Mathematics
\textbf{91}, Springer, New York, 1983. doi:10.1007/978-1-4612-1146-4.
\bibitem{CartanDieudonne} E.~Artin, \emph{Geometric Algebra}, Interscience, New York, 1957, Chapter~III, Theorem~3.20
(p.~129), the Cartan--Dieudonn\'e theorem: every isometry of a non-singular $n$-dimensional quadratic
space is a product of at most $n$ reflections. For the Euclidean case see also J.~Gallier,
\emph{Geometric Methods and Applications}, 2nd ed., Springer, 2011, Theorem~8.1.
\bibitem{Wilker} J.~B.~Wilker, \emph{The quaternion formalism for M\"obius groups in four or fewer
dimensions}, Linear Algebra Appl. \textbf{190} (1993), 99--136. doi:10.1016/0024-3795(93)90222-A.
\bibitem{Vahlen} K.~Th.~Vahlen, \emph{Ueber Bewegungen und complexe Zahlen}, Math. Ann. \textbf{55} (1902),
no.~4, 585--593. doi:10.1007/BF01450354. (Context only: our normal form does not use the Vahlen
characterization.)
\bibitem{Ahlfors85} L.~V.~Ahlfors, \emph{M\"obius transformations and Clifford numbers}, in \emph{Differential
Geometry and Complex Analysis} (I.~Chavel, H.~M.~Farkas, eds.), Springer, Berlin, 1985, 65--73.
doi:10.1007/978-3-642-69828-6\_5. (Context only.)
\bibitem{EGM} J.~Elstrodt, F.~Grunewald, J.~Mennicke, \emph{Vahlen's group of Clifford matrices and
spin-groups}, Math.~Z. \textbf{196} (1987), no.~3, 369--390. doi:10.1007/BF01200359. (Context only.)
\bibitem{Waterman} P.~L.~Waterman, \emph{M\"obius transformations in several dimensions}, Adv. Math.
\textbf{101} (1993), no.~1, 87--113. doi:10.1006/aima.1993.1043. (Context only.)
\bibitem{Waldschmidt} M.~Waldschmidt, \emph{Open Diophantine problems}, Mosc. Math. J. \textbf{4}
(2004), no.~1, 245--305; doi:10.17323/1609-4514-2004-4-1-245-305; arXiv:math/0312440.
\bibitem{Baker} A.~Baker, \emph{Transcendental Number Theory}, Cambridge University Press, 1975,
Chapter~2 (the Gelfond--Schneider theorem).
\bibitem{Falconer} K.~J.~Falconer, \emph{The realization of distances in measurable subsets covering $\R^n$},
J. Combin. Theory Ser. A \textbf{31} (1981), no.~2, 184--189. doi:10.1016/0097-3165(81)90014-5.
\bibitem{ShelahSoifer} S.~Shelah, A.~Soifer, \emph{Axiom of choice and chromatic number of the plane}, J. Combin.
Theory Ser. A \textbf{103} (2003), no.~2, 387--391. doi:10.1016/S0097-3165(03)00102-X.
\bibitem{Payne} M.~S.~Payne, \emph{Unit distance graphs with ambiguous chromatic number}, Electron. J. Combin.
\textbf{16} (2009), no.~1, Note~31. doi:10.37236/269.
\bibitem{Bourgain} J.~Bourgain, \emph{A Szemer\'edi type theorem for sets of positive density in $\R^k$}, Israel
J. Math. \textbf{54} (1986), no.~3, 307--316. doi:10.1007/BF02764959.
\bibitem{CastroSilva} D.~Castro-Silva, F.~M.~de Oliveira Filho, L.~Slot, F.~Vallentin, \emph{A recursive Lov\'asz
theta number for simplex-avoiding sets}, Proc. Amer. Math. Soc. \textbf{150} (2022), no.~8,
3307--3322. doi:10.1090/proc/15940. arXiv:2106.09360.
\bibitem{GuruswamiLi} V.~Guruswami, S.~Li, \emph{Density Frankl--R\"odl on the sphere}, arXiv:2505.09839 (2025).
\bibitem{Sun} W.~Sun, \emph{Spherical higher order Fourier analysis over finite fields IV: an application to the
Geometric Ramsey Conjecture}, arXiv:2312.06649 (2023).
\bibitem{Zapletal} J.~Zapletal, \emph{Coloring equilateral triangles}, arXiv:2503.18295 (2025).
\bibitem{Behague} N.~Behague, \emph{Nearly all known Euclidean Ramsey sets are subsoluble},
arXiv:2510.15677 (2025).
\bibitem{Moore} K.~Moore, \emph{A pyramid with a Ramsey base is Ramsey}, arXiv:2608.09649 (2026).
\bibitem{Mirabi} M.~Mirabi, \emph{One-point extensions of Euclidean Ramsey sets}, arXiv:2608.11736 (2026).
\bibitem{DFPP} G.~Dam\'asdi, N.~Frankl, J.~Pach, D.~P\'alv\"olgyi, \emph{Monochromatic configurations
on a circle}, arXiv:2504.10687 (2025).
\bibitem{BukhKarasev} B.~Bukh, R.~Karasev, \emph{Suborbits in Knaster's problem}, Bull. Lond. Math. Soc.
\textbf{46} (2014), no.~2, 269--278; arXiv:1212.5351.
\bibitem{DolnikovKarasev} V.~L.~Dol'nikov, R.~N.~Karasev, \emph{Dvoretzky type theorems for multivariate
polynomials and sections of convex bodies}, Geom. Funct. Anal. \textbf{21} (2011), no.~2, 301--318.
doi:10.1007/s00039-011-0109-8. arXiv:1009.0392.
\bibitem{Gromov} M.~Gromov, \emph{Number of questions}, preprint, 26 April 2014, \S1.12;
\url{https://www.ihes.fr/~gromov/recent/}.
\bibitem{vdDries} L.~van den Dries, \emph{Tame Topology and O-minimal Structures}, London Math. Soc.
Lecture Note Ser. \textbf{248}, Cambridge University Press, 1998.
\bibitem{Comparator} leanprover/comparator: stock Comparator built at tag v4.34.0 and rebuilt with Lean
v4.34.1, run under Landlock (ropes AV1--AV5); and \texttt{lake comparator} as bundled in
\texttt{leanprover/lean4:v4.35.0-rc2}, with the NanoDa and con-ron kernels (the package \cite{MR}).
\bibitem{vibemathed} VibeMathed, a catalogue of AI-produced mathematical results, \url{https://vibemathed.com/}
(data at \url{https://github.com/mrconter1/vibemathed}).
\end{thebibliography}

\end{document}
